\chapter{摩西五经总结}

\section{出埃及，民数记，申命记时间的顺序}
\begin{figure}[H]
\centering
\includegraphics[width=7.in]{ZongJie/figs/AMoXiWuJing/AZongJie1}
\caption{亚述}
\label{fig:sub1}
\end{figure}


\begin{figure} 
\centering
\includegraphics[width=6.in]{ZongJie/figs/AMoXiWuJing/AZongJie2}
\caption{亚述}
\label{fig:sub1}
\end{figure}


《摩西五经》相当于五个手指，核心的手指在中间，中指就相当于《利未记》，因为《利未记》里面有一个核心，就是如何献祭的一些制度。我们人一生中最大的问题就是如何把罪除掉，如何向上帝献祭，使我们罪得赦免，人又该遵循哪一些规则，这就成为整个救赎计划核心的部分。中指之外还有另外前面的两个指头，大拇指就是《创世记》，食指就是《出埃及记》，无名指就是《民数记》，它是《出埃及记》的一个延续，是迈向迦南地的一个旅程的继续。最后就是小拇指，重新叙述了上帝与以色列人所订立的规范、法则，这就是《申命记》。


\begin{figure} 
\centering
\includegraphics[width=6.in]{ZongJie/figs/AMoXiWuJing/anotherRoute}
\caption{\href{http://www.biblelib.ca/2021/01/page/5/}{圣经书院：真正的西奈旷野} }
\label{fig:sub1}
\end{figure}

\begin{figure} 
\centering
\includegraphics[width=6.in]{ZongJie/figs/AMoXiWuJing/newRoute}
\caption{\href{https://posts.careerengine.us/p/5f0a7c84e1bc557c6e745808L}{真正的西奈旷野} }
\label{fig:sub1}
\end{figure}


\begin{figure} 
\centering
\includegraphics[width=6.in]{ZongJie/figs/AMoXiWuJing/newRoute2}
\caption{\href{http://bible4tyro.blogspot.com/2017/10/blog-post_17.html}{小呆羊读圣经：真正的西奈旷野} }
\label{fig:sub1}
\end{figure}

\begin{table}[H]
\centering
\begin{tabular}{p{2cm}|p{11cm}|p{3.5cm}}\hline\hline
\multicolumn{1}{c|}{时间/地点}&\multicolumn{1}{c|}{事件/人物}&\multicolumn{1}{c}{圣经章节}\\\hline
1/15蘭塞&以色列人被赶出埃及&出12:37,40,41/民33:5\\\hline
疏割&第一个营地&出13:20\\\hline
以倘&在曠野邊的以倘安營,日間雲柱，夜間火柱總不離開百姓的面前&出13:20-22\\\hline
比哈希錄&安營在比哈希錄前、密奪和海的中間，對著巴力洗分，靠近海邊安營&出14:2\\\hline
紅海邊&以色列人走乾地,海水淹沒埃及車輛和馬兵&出14:1-31\\\hline
三天書珥&從紅海往前行到了書珥的曠野找不著水&出15:22\\\hline
瑪拉&摩西把樹丟在水裡變苦水為甜&出15:23-26\\\hline
以琳&十二股水泉七十棵棕樹邊安營&出15:27\\\hline
2/15汎曠野&耶和華賜以色列人嗎哪和鵪鶉為食物&出16:1,13,14\\\hline
利非訂&沒有水喝爭鬧,試探耶和華,那地方叫瑪撒又叫米利巴&出17:1,7\\\hline
利非訂&與亞瑪力人爭戰&出17:8-16\\\hline
利非訂&米甸祭司葉忒羅建议揀選千/百/五十/十夫長管理百姓&出18:1-27/申1:9-18\\\hline

4/15西乃山&來到西奈的曠野&出19:1\\\hline
4/18&耶和華在眾百姓眼前降臨在西奈山上&出19:16-25\\\hline
4/18&耶和華曉諭摩西十戒&出20:1-17\\\hline
4/18&二次曉諭摩西築壇/待僕/殺人/律法&出20:18-23:33\\\hline
4/19清早起&摩西將命令寫上,築壇立十二根柱子灑血灑立約&出24:1-8\\\hline
雲蓋山六天&摩西/約書亞/亞倫/拿答/亞比戶/七十長老上山&出24:15-16\\\hline
4/25第七天&四十晝夜耶和華曉諭摩西製造帳幕设立祭司，給兩塊法版&出24:17-31:18\\\hline
山下&亞倫領百姓鑄一隻金牛犢&出32:1-6\\\hline
山上&摩西代求&出32:7-14\\\hline
山下&怒摔法版，利未的子孫照摩西話殺三千百姓&出32:15-29\\\hline
6/5第二天&摩西回耶和華那裡代求，在營外會幕&出32:30-34:4\\\hline
6/6清晨&鑿出兩塊石版上西奈山四十晝夜(6/6-7/15)不吃飯不喝水&出34:5-28\\\hline
第二年1/1&建會幕,立起帳幕完工&出35-40\\\hline
\end{tabular}
\caption{出埃及記}
\end{table}



\begin{table}[H]
\centering
\begin{tabular}{p{3.2cm}|p{11cm}|p{3.2cm}}\hline\hline
\multicolumn{1}{c|}{时间/地点}&\multicolumn{1}{c|}{事件/人物}&\multicolumn{1}{c}{圣经章节}\\\hline
第二年2/1&在會幕中曉諭摩西核男丁總數,各歸各纛&民1-2\\\hline
&分別利未為頭生的歸耶和華為聖,分配利未各族職任&民3-4\\\hline
&潔淨營地，疑妻行淫之法，拿細耳人之例&民5-6\\\hline
&會幕建立族長獻禮&民7\\\hline
第二年1/14,2/14&潔淨利未人之例,守逾越節,製銀號&民8:1-10:10\\\hline
第二年2/20&離開西奈曠野，行三天路程在巴蘭的曠野&民10:11-36\\\hline
他備拉&百姓發怨言火燒在他們中間&民11:1-3\\\hline
基博羅哈他瓦&閒雜人起貪心,百姓哭號,立七十長老,捕取鵪鶉,重災殃擊殺百姓&民11:4-34\\\hline
哈洗錄&米利暗和亞倫因古實女子毀謗摩西，米利暗長大痲瘋關營外七天&民11:35-12:15\\\hline
%利提瑪&&\\\hline
%臨門帕烈&&\\\hline
%立拿&&\\\hline
%勒撒&&\\\hline
%基希拉他&&\\\hline
%沙斐山&&\\\hline
%哈拉大&&\\\hline
%瑪吉希錄&&\\\hline
%他哈&&\\\hline
%他拉&&\\\hline
%密加&&\\\hline
%哈摩拿&&\\\hline
%摩西錄&&\\\hline
%比尼亞干&&\\\hline
%曷哈及甲&&\\\hline
%約巴他&&\\\hline
%阿博拿&&\\\hline
%以旬迦別&&\\\hline
%尋的曠野加低斯&&\\\hline


加低斯巴尼亞&何烈山到巴蘭曠野11天的路,窺探迦南地40天；以色列眾人發怨言&民12-14/申1\\\hline
\textcolor{blue}{西珥山繞行38年}&可拉黨 250 首領叛逆叛亂，民遭疫癘，亞倫之杖發芽開花結熟杏&民16-17\\\hline
\textcolor{blue}{西珥山繞行38年}&祭司與利未人之職和酬勞，除汙穢的水&民18-19\\\hline
正月加低斯&\textcolor{red}{米利暗死}在尋的曠野加低斯&民20:1\\\hline
米利巴&摩西因无水杖擊打磐石受牵累不得領會眾進迦南&民20:2-13\\\hline
加低斯&摩西差遣使者去見以東王&民20:14-21/申2\\\hline

40年5/1何珥山&\textcolor{red}{亞倫死}在何珥山&民20:22-29\\\hline
何珥瑪&迦南人亞拉得王被毀&民21:1-3\\\hline
何珥山繞以東&製造銅火蛇&民21:5-9\\\hline
%比珥&&\\\hline
阿伯&&民21:10\\\hline
以耶亞巴琳&與摩押相對的曠野，向日出之地&民21:11\\\hline
撒烈谷&&民21:12\\\hline
亞嫩河&這亞嫩河是在曠野，從亞摩利的境界流出來的，是摩押的邊界&民21:13-15\\\hline
比珥&從前耶和華吩咐摩西說：「招聚百姓，我好給他們水喝」，說的就是這井。&民21:16-18\\\hline
瑪他拿&&民21:18\\\hline
拿哈列&&民21:19\\\hline

巴末&&民21:19-20\\\hline
摩押地的谷&&民21:20\\\hline
毗斯迦山頂&&民21:20\\\hline
亞嫩谷&以色列人擊敗亞摩利人希實本王西宏&民21:21-32/申2\\\hline
以得來&巴珊王噩和他的眾民都出來與以色列人交戰&民21:33-35/申3\\\hline
摩押平原&摩押王召巴蘭咒詛以色列人&民22-25\\\hline
&第二次數點以色列人&民26:1-65\\\hline
&瑪拿西族西羅非哈的女兒們求產業&民27:1-11\\\hline
&立嫩的兒子約書亞&民27:12-23/申3:21\\\hline
&與米甸人打仗&民31\\\hline
&河東的兩個半支派求業於約旦河東&民32/申3:12-20\\\hline
40年11/1摩押平原&第一篇讲论:回顾历史背景。約但河邊耶利哥對面&申1:1-4:40\\\hline
&設立河東之三座逃城&申4:41-43\\\hline
&第二篇讲论:基本条约申5-11/特殊条约(细节)申12-26&申5-26\\\hline
&第三篇讲论:祝福与咒诅&申27-28\\\hline
&第四篇讲论:圣约被毁更新&申29-30\\\hline
&第五篇讲论:预言以色列不忠,祝福以色列民&申31-34\\\hline
尼波山&\textcolor{red}{摩西從摩押平原登尼波山,120歲死在摩押}&申34:6\\\hline
&以色列人在摩押平原為摩西哀哭了三十日&申34:6\\\hline
&嫩的兒子約書亞接替摩西&申34:9-11\\\hline
41年1/10&百姓從約旦河裡上來，就在吉甲&書4:19\\\hline
\end{tabular}
\caption{民數記、申命記}
\end{table}



\section{創世記里的基督}

\subsection{安息}
\begin{table}[H]
\centering
\begin{tabular}{p{2cm}|p{14.5cm}}\hline\hline
\multicolumn{1}{c|}{圣经章节}&\multicolumn{1}{c}{事件/人物}\\\hline
創 2:1-3&天地萬物都造齊了。\textcolor{red}{到第七日，神造物的工已經完畢，就在第七日歇了他一切的工，安息了}。神賜福給第七日，定為聖日，因為在這日神歇了他一切創造的工，就安息了。\\\hline

出16:23&
摩西對他們說：「耶和華這樣說：『明天是聖安息日，是向耶和華守的聖安息日。你們要烤的就烤了，要煮的就煮了，所剩下的都留到早晨。』」\\\hline

出 20:8-11&当记念安息日，守为圣日。六日要劳碌做你一切的工，但第七日是向耶和华你神当守的安息日。这一日你和你的儿女、仆婢、牲畜，并你城里寄居的客旅，无论何工都不可做。因为六日之内，耶和华造天、地、海和其中的万物，第七日便安息，所以\textcolor{red}{耶和华赐福于安息日，定为圣日}。\\\hline

出 23:10-12&六年你要耕種田地，收藏土產， 11 只是第七年要叫地歇息，不耕不種，使你民中的窮人有吃的，他們所剩下的，野獸可以吃。你的葡萄園和橄欖園也要照樣辦理。 12 六日你要做工，第七日要安息，使牛、驢可以歇息，並使你婢女的兒子和寄居的都可以舒暢。 13 凡我對你們說的話，你們要謹守。別神的名，你不可提，也不可從你口中傳說。\\\hline

出 31:12-17&耶和華曉諭摩西說： 13 「你要吩咐以色列人說：『你們務要守我的安息日，因為這是\textcolor{red}{你我之間世世代代的證據，使你們知道我耶和華是叫你們成為聖的}。 14 所以你們要守安息日，以為聖日。凡干犯這日的，必要把他治死；凡在這日做工的，必從民中剪除。 15 六日要做工，但第七日是安息聖日，是向耶和華守為聖的。凡在安息日做工的，必要把他治死。 16 故此，以色列人要世世代代守安息日為永遠的約。 17 這是我和以色列人永遠的證據，因為六日之內耶和華造天地，第七日便安息舒暢。』」\\\hline
出 34:21&你六日要做工，第七日要安息；雖在耕種收割的時候，也要安息。\\\hline

出 35:1-3&
摩西招聚以色列全會眾，對他們說：「這是耶和華所吩咐的話，叫你們照著行。 2 六日要做工，第七日乃為聖日，當向耶和華守為安息聖日。凡這日之內做工的，必把他治死。 3 當安息日，不可在你們一切的住處生火。」\\\hline

利 19:3&你們各人都當孝敬父母，也要守我的安息日。我是耶和華你們的神。\\\hline


利 19:30&你們要守我的安息日，敬我的聖所。我是耶和華。\\\hline

利 23:3&
「六日要做工，第七日是聖安息日，當有聖會，你們什麼工都不可做。這是在你們一切的住處向耶和華守的安息日。\\\hline

利未記 25:2&
「你曉諭以色列人說：你們到了我所賜你們那地的時候，地就要向耶和華守安息。\\\hline
利 26:2&你們要守我的安息日，敬我的聖所。我是耶和華。\\\hline
申 5:12-15&當照耶和華你神所吩咐的守安息日為聖日。 13 六日要勞碌做你一切的工， 14 但第七日是向耶和華你神當守的安息日。這一日，你和你的兒女、僕婢、牛、驢、牲畜，並在你城裡寄居的客旅，無論何工都不可做，使你的僕婢可以和你一樣安息。 15 你也要記念你在埃及地做過奴僕，耶和華你神用大能的手和伸出來的膀臂將你從那裡領出來，因此耶和華你的神吩咐你守安息日。\\\hline



賽 56:1-8&
耶和華如此說：「你們當守公平，行公義，因我的救恩臨近，我的公義將要顯現。 2 謹守安息日而不干犯，禁止己手而不作惡，如此行、如此持守的人，便為有福！」 3 與耶和華聯合的外邦人不要說：「耶和華必定將我從他民中分別出來。」太監也不要說：「我是枯樹。」 4 因為耶和華如此說：「那些謹守我的安息日，揀選我所喜悅的事，持守我約的太監， 5 我必使他們在我殿中、在我牆內有紀念，有名號，比有兒女的更美。我必賜他們永遠的名，不能剪除。還有那些與耶和華聯合的外邦人——要侍奉他，要愛耶和華的名，要做他的僕人——就是凡守安息日不干犯，又持守他約的人， 7 我必領他們到我的聖山，使他們在禱告我的殿中喜樂。他們的燔祭和平安祭在我壇上必蒙悅納，因我的殿必稱為萬民禱告的殿。」 8 主耶和華，就是招聚以色列被趕散的，說：「在這被招聚的人以外，我還要招聚別人歸併他們。\\\hline

賽 58:13-14&你若在安息日掉轉你的腳步，在我聖日不以操作為喜樂，稱安息日為可喜樂的，稱耶和華的聖日為可尊重的，而且尊敬這日，不辦自己的私事，不隨自己的私意，不說自己的私話，你就以耶和華為樂，耶和華要使你乘駕地的高處，又以你祖雅各的產業養育你。」這是耶和華親口說的。\\\hline
\end{tabular}
\caption{安息日1}
\end{table}

\begin{table}[H]
\centering
\begin{tabular}{p{2cm}|p{14.5cm}}\hline\hline
\multicolumn{1}{c|}{圣经章节}&\multicolumn{1}{c}{事件/人物}\\\hline
耶 17:21-27&耶和華如此說：你們要謹慎，不要在安息日擔什麼擔子進入耶路撒冷的各門， 22 也不要在安息日從家中擔出擔子去。無論何工都不可做，只要以安息日為聖日，正如我所吩咐你們列祖的。 23 他們卻不聽從，不側耳而聽，竟硬著頸項不聽，不受教訓。
24 『耶和華說：你們若留意聽從我，在安息日不擔什麼擔子進入這城的各門，只以安息日為聖日，在那日無論何工都不做， 25 那時就有坐大衛寶座的君王和首領，他們與猶大人並耶路撒冷的居民，或坐車或騎馬進入這城的各門，而且這城必存到永遠。 26 也必有人從猶大城邑和耶路撒冷四圍的各處，從便雅憫地、高原、山地並南地而來，都帶燔祭、平安祭、素祭和乳香並感謝祭到耶和華的殿去。 27 你們若不聽從我，不以安息日為聖日，仍在安息日擔擔子進入耶路撒冷的各門，我必在各門中點火，這火也必燒毀耶路撒冷的宮殿，不能熄滅。』\\\hline

結 20:12&
又將我的安息日賜給他們，好在我與他們中間為證據，使他們知道我耶和華是叫他們成為聖的。\\\hline

結 20:20&
且以我的安息日為聖。這日在我與你們中間為證據，使你們知道我是耶和華你們的神。」\\\hline

結 37:28&
我的聖所在以色列人中間直到永遠，外邦人就必知道我是叫以色列成為聖的耶和華。』」\\\hline

結 44:24&
有爭訟的事，他們應當站立判斷，要按我的典章判斷。在我一切的節期，必守我的律法、條例，也必以我的安息日為聖日。\\\hline


尼 9:14&
又使他們知道你的安息聖日，並藉你僕人摩西傳給他們誡命、條例、律法。\\\hline

尼 13:15&
那些日子，我在猶大見有人在安息日榨酒，搬運禾捆馱在驢上，又把酒、葡萄、無花果和各樣的擔子在安息日擔入耶路撒冷，我就在他們賣食物的那日警戒他們。\\\hline


太11:25-30&那時，耶穌說：「父啊，天地的主，我感謝你！因為你將這些事向聰明通達人就藏起來，向嬰孩就顯出來。 26 父啊，是的，因為你的美意本是如此。 27 一切所有的，都是我父交付我的。除了父，沒有人知道子；除了子和子所願意指示的，沒有人知道父。 28 凡勞苦擔重擔的人，可以到我這裡來，我就使你們得安息。 29 我心裡柔和謙卑，你們當負我的軛，學我的樣式，這樣你們心裡就必得享安息。 30 因為我的軛是容易的，我的擔子是輕省的。」\\\hline


太12:1-8&那時，耶穌在安息日從麥地經過。他的門徒餓了，就掐起麥穗來吃。法利賽人看見，就對耶穌說：「看哪，你的門徒做安息日不可做的事了。」耶穌對他們說：「經上記著大衛和跟從他的人飢餓之時所做的事，你們沒有念過嗎？他怎麼進了神的殿，吃了陳設餅？這餅不是他和跟從他的人可以吃的，唯獨祭司才可以吃。再者，律法上所記的，當安息日祭司在殿裡犯了安息日還是沒有罪，你們沒有念過嗎？但我告訴你們：在這裡有一人比殿更大。『我喜愛憐恤，不喜愛祭祀』，你們若明白這話的意思，就不將無罪的當做有罪的了。因為人子是安息日的主。」\\\hline


太12:9-13&耶穌離開那地方，進了一個會堂。 10 那裡有一個人枯乾了一隻手。有人問耶穌說：「安息日治病可以不可以？」意思是要控告他。 11 耶穌說：「你們中間誰有一隻羊當安息日掉在坑裡，不把牠抓住拉上來呢？ 12 人比羊何等貴重呢！所以，在安息日做善事是可以的。」 13 於是對那人說：「伸出手來！」他把手一伸，手就復了原，和那隻手一樣。\\\hline


可 2:27&又對他們說：「安息日是為人設立的，人不是為安息日設立的。\\\hline

路 13:14&
管會堂的因為耶穌在安息日治病，就氣憤憤地對眾人說：「有六日應當做工，那六日之內可以來求醫，在安息日卻不可。」\\\hline
約 5:17&耶穌就對他們說：「我父做事直到如今，我也做事。」\\\hline
%希伯來書 4:4&論到第七日，有一處說：「到第七日神就歇了他一切的工。」\\\hline

來 4:3-11&但我們已經相信的人得以進入那安息，正如神所說：「我在怒中起誓說：『他們斷不可進入我的安息。』」其實造物之工，從創世以來已經成全了。 4 論到第七日，有一處說：「到第七日神就歇了他一切的工。」 5 又有一處說：「他們斷不可進入我的安息。」 6 既有必進安息的人，那先前聽見福音的因為不信從，不得進去， 7 所以過了多年，就在大衛的書上又限定一日，如以上所引的說：「你們今日若聽他的話，就不可硬著心。」 8 若是約書亞已叫他們享了安息，後來神就不再提別的日子了。 9 這樣看來，必另有一安息日的安息為神的子民存留。 10 因為那進入安息的，乃是歇了自己的工，正如神歇了他的工一樣。 11 所以，我們務必竭力進入那安息，免得有人學那不信從的樣子跌倒了。\\\hline


\end{tabular}
\caption{安息日}
\end{table}




\subsection{成聖}
\begin{table}[H]
\centering
\begin{tabular}{p{3.cm}|p{14cm}}\hline\hline
\multicolumn{1}{c|}{圣经章节}&\multicolumn{1}{c}{事件/人物}\\\hline
利未記 10:3&
於是摩西對亞倫說：「這就是耶和華所說：『我在親近我的人中要顯為聖，在眾民面前我要得榮耀。』」亞倫就默默不言。\\\hline

利 20:8&
你們要謹守遵行我的律例，我是叫你們成聖的耶和華。\\\hline

利21:8&
所以你要使他成聖，因為他奉獻你神的食物；你要以他為聖，因為我，使你們成聖的耶和華是聖的。\\\hline

結 37:28&
我的聖所在以色列人中間直到永遠，外邦人就必知道我是叫以色列成為聖的耶和華。』」\\\hline
約17:17,19&
求你用真理使他們成聖，你的道就是真理。\\\hline

帖前 5:23&
願賜平安的神親自使你們全然成聖！又願你們的靈與魂與身子得蒙保守，在我們主耶穌基督降臨的時候完全無可指摘！
\\\hline
\end{tabular}
\caption{民數記、申命記}
\end{table}


\section{出埃及里的基督}
\subsection{摩西出生法老击杀男婴预表耶稣出生希律击杀男婴}
《出埃及记》第一章，有一个法老发出的要击杀男婴的残酷指令，这奠定了一个环境，摩西就是在此环境下出生的。这个击杀男婴的指令我们看到摩西出生的背景，实际上这就是预表了耶稣出生的情景。

《马太福音》第二章中写到，有几位东方博士来到耶路撒冷朝拜新生王，当时是希律做王。他觉得这里有一个新的王要来跟他竞争，就下命令要把这个男婴杀害。所以他吩咐博士说：你们找到他了，回头告诉我。可是主的使者指示这些博士，最后他们没有回去。

《马太福音》2：12 -18“博士因为在梦中被主指示不要回去见希律，就从别的路回本地去了。他们去后，有主的使者向约瑟梦中显现，说：“起来！带着小孩子同他母亲逃往埃及，住在那里，等我吩咐你；因为希律必寻找小孩子，要除灭他。”约瑟就起来，夜间带着小孩子和他母亲往埃及去，住在那里，直到希律死了。这是要应验主藉先知所说的话，说：“我从埃及召出我的儿子来。”希律见自己被博士愚弄，就大大发怒，差人将伯利恒城里并四境所有的男孩，照着他向博士仔细查问的时候，凡两岁以里的，都杀尽了。这就应了先知耶利米的话，说：‘在拉玛听见号啕大哭的声音，是拉结哭她儿女，不肯受安慰，因为他们都不在了’”。

希律王所发的这个命令就和《出埃及记》第一章所说的相似。埃及王对两位收生婆说：“你们为希伯来妇人收生，看她们临盆的时候，若是男孩，就把他杀了；若是女孩，就留她存活。”但是收生婆敬畏上帝，不照埃及王的吩咐行，竟存留男孩的性命。（出埃及记1:16-17）

摩西就在这样一个大的背景下出生。他出生的时候，因为他的父母敬畏上帝，不愿意把他丢弃，把他隐藏了三个月，后来实在是藏不住了，就把他放在一个蒲草箱子里，涂上石漆和石油，把孩子放在里面，然后又把他放在河边，让他的姐姐远远的在那里看着，看看会发生什么事情。（出1:1-4）

%\subsection{摩西出生被法老的女儿领养预表耶稣由童女所生}
%后来摩西就被法老的女儿领养。这个女孩当时并没有结婚，换句话说，她是一个童女，而摩西是被一个童女所领养的。从这里我们看到摩西也预表了耶稣，是要由童女所生，童女马利亚怀孕生下耶稣。

%《出埃及记》的第一章、第二章，摩西出生的背景，就向我们展示了耶稣出生的背景。这是第一个问题。

\subsection{使者从荆棘的火焰中向摩西显现预表耶稣荆棘作冠冕}
看见上帝在荆棘丛中。有一天，摩西牧养他岳父米甸祭司叶忒罗的羊群，来到何烈山。耶和华的使者从荆棘的火焰中向摩西显现。圣经在《出埃及记》中记载说，”摩西观看，不料，荆棘被火烧着，却没有烧毁。摩西说：“我要过去看这大异象，这荆棘为何没有烧坏呢？”耶和华上帝见他过去要看，就从荆棘里呼叫说：“摩西！摩西！”他说：“我在这里。”上帝说：“不要近前来。当把你脚上的鞋脱下来，因为你所站之地是圣地。”又说：“我是你父亲的上帝，是亚伯拉罕的上帝，以撒的上帝，雅各的上帝。”摩西蒙上脸，因为怕看上帝。”（出3：1-6）

上帝所显现的地点是在何烈山,在荆棘中显现。这个荆棘，我们看到在耶稣受审判的时候，那恶者用荆棘作冠冕带在他的头上。荆棘在万物中是很卑微的一种植物。慈悲的上帝用最卑微之物遮掩了他的荣耀，使摩西能以看见而仍得存活。

\subsection{从埃及地出来带的金器、银器预表教会身上所披戴的荣耀的身体}
我们知道，当人犯罪被上帝从伊甸园赶出来的时候，他们是以什么样的方式出来的呢？我们的始祖亚当夏娃犯罪以后，他们身上荣耀的衣服就离开他们了，他们身上所披的就是自己用无花果树的叶子做的衣服。出到伊甸园之外，上帝用皮子给他们做了衣服。

而这里说，你们从埃及地出来的时候，你们儿女身上所穿戴的都是金器、银器和衣服，就不是无花果树叶的那种衣服了。上帝在这里预表，将来耶稣带领他的教会回到天家的时候，他的教会身上所披戴的，不再是无花果树的叶子，而是复活以后荣耀的身体。就像被创造的时候那样，有着荣耀、尊贵为冠冕。


\subsection{逾越节}
\subsubsection{无酵饼}
逾越节的有关记述在诸多方面预表了基督是为我们赎罪牺牲的。耶稣为逾越节带来的一个新的制度。出12:37,“以色列人从兰塞起行，往疏割去，除了妇人孩子，步行的男人约有六十万。”历经400年，以色列人从七十个人，发展到六十万，再加上妇女和孩子，保守的说法也应该有200万之众。而且出12:38-39“又有许多闲杂人，并有羊群牛群，和他们一同上去。他们用埃及带出来的生面烤成无酵饼，这生面原没有发起，因为他们被催逼离开埃及不能耽延，也没有为自己预备什么食物。”这是他们出离埃及的一个情形。



\subsubsection{初十日取羊羔}
以色列人出埃及的那一天，就是逾越节。逾越节是在第十二章的一开篇。耶和华在埃及地晓谕摩西、亚伦说：“你们要以本月为正月，为一年之首。你们吩咐以色列全会众说：本月初十日，各人要按着父家取羊羔，一家一只。若是一家的人太少，吃不了一只羊羔，本人就要和他隔壁的邻舍共取一只。你们预备羊羔，要按着人数和饭量计算。要无残疾、一岁的公羊羔，你们或从绵羊里取，或从山羊里取，都可以。要留到本月十四日，在黄昏的时候，以色列全会众把羊羔宰了。各家要取点血，涂在吃羊羔的房屋左右的门框上和门楣上。当夜要吃羊羔的肉，用火烤了，与无酵饼和苦菜同吃。（出12:1-8）

\subsubsection{头生}
出12：12“因为那夜我要巡行埃及地，把埃及地一切头生的，无论是人是牲畜，都击杀了，又要败坏埃及一切的神。我是耶和华。”这里我们看到，上帝要败坏埃及地一切的神。耶稣在小的时候，他其实也是住在埃及地的。耶稣就是头生的儿子，也可以说是埃及地头生的儿子。同时以动物牺牲的形式来说，他是头生的羔羊、头生的公牛犊等等，都是来指代耶稣。


\subsubsection{血要在房屋上作记号}
出12：13-14“这血要在你们所住的房屋上作记号，我一见这血，就越过你们去。我击杀埃及地头生的时候，灾殃必不临到你们身上灭你们。你们要记念这日，守为耶和华的节，作为你们世世代代永远的定例。”这里我们看到，《出埃及记》所制定的逾越节，这些细节，在耶稣的身上都有详细的应验。这个逾越节的羔羊是没有瑕疵的。在新约里面就说，“看哪，上帝的羔羊，除去世人罪孽的。”（约1：29）他是没有瑕疵的，他就是上帝所预备的逾越节的羔羊，为要背负世人的罪而舍去，牺牲自己的生命。到学习新约的时候，我们还会回到这个主题。


\subsubsection{羊羔的骨头一根也不可折断}
出12:46“应当在一个房子里吃，不可把一点肉从房子里带到外头去。羊羔的骨头一根也不可折断。”很明显，特别是最后这句话“羊羔的骨头一根也不可折断”，在耶稣的身上是得到了应验的。约19:31-36是讲耶稣被钉十字架之后的情形：“犹太人因这日是预备日，又因那安息日是个大日，就求彼拉多叫人打断他们的腿，把他们拿去，免得尸首当安息日留在十字架上。于是兵丁来，把头一个人的腿，并与耶稣同钉第二个人的腿都打断了。只是来到耶稣那里，见他已经死了，就不打断他的腿。惟有一个兵拿枪扎他的肋旁，随即有血和水流出来。看见这事的那人就作见证，他的见证也是真的，并且他知道自己所说的是真的，叫你们也可以信。这些事成了，为要应验经上的话说：‘他的骨头一根也不可折断。’”这里我们看到，以色列人当初被吩咐吃全羊的时候，羊的骨头是不可以折断的，并只能在一个房子里吃，不可到处走。耶稣是在营外，在以色列耶路撒冷的城外被钉十字架的。

\subsection{日间有云柱、夜间有火柱带领}

第十三章是上帝教以色列人如何守除酵节；以色列人在旷野行走时，日间有云柱、夜间有火柱带领，并总不离开百姓的面前。


\subsection{摩西诗篇}
出15:2“耶和华是我的力量、我的诗歌，也成了我的拯救。这是我的上帝，我要赞美他；是我父亲的上帝，我要尊崇他。”这里所说的力量、诗歌、拯救者就是指耶稣

摩西在这个诗篇中给我们留下了很多很宝贵的赞美上帝的经文；出15:11-13“耶和华啊，众神之中谁能像你?谁能像你至圣至荣，可颂可畏，施行奇事？你伸出右手，地便吞灭他们。你凭慈爱领了你所赎的百姓；你凭能力领他们到了你的圣所。”这节经文非常得重要。这时，他们刚刚过了红海，还没有造圣所，也没有到达迦南地，但是摩西在诗篇里就说：你凭慈爱领了你所赎的百姓，你凭能力领他们到了你的圣所。

出15:17“你要将他们领进去，栽于你产业的山上。耶和华啊，就是你为自己所造的住处；主啊，就是你手所建立的圣所。”在这诗篇里，摩西两次提到了“圣所”，将来我们在《遍览圣经觅圣所》的时候，我们会再次回到这两节经文，现在只是提出，请大家注意。我们以前提到圣所一般都是引用出25:8“又当为我造圣所。”但是“圣所”这个词第一次出现是在《出埃及记》第十五章

\subsection{玛拉摩西把树丢在水里，水就变甜了}
第十五章还谈到了另外一件事情。到了玛拉这个地方，里边的水不能喝，因为水苦，所以就把地名叫作玛拉。出15:24-25“百姓就向摩西发怨言，说：‘我们喝什么呢？’摩西呼求耶和华，耶和华指示他一颗树，他把树丢在水里，水就变甜了。”

《圣经》里常常用树来形容义人。比方说，在诗1：1-3节，就这样说：“不从恶人的计谋，不站罪人的道路，不坐亵慢人的座位，惟喜爱耶和华的律法，昼夜思想，这人便为有福。他要像一颗树栽在溪水旁，按时结果子，叶子也不枯干。凡他所作的尽都顺利。”

义人是这样的，那么谁是义者呢？经上说，耶稣就是那个义者。另外，义人作为一棵树的比喻用在耶稣身上，还有约15:1耶稣说：“我是真葡萄树，我父是栽培的人。”同时他也希望我们都能连在这棵树上多结果子。所以，耶稣把自己比喻成一棵树在新约中也应验了。当摩西把树丢在水中，水就变甜了，同样耶稣使我们在罪中的苦难因为他，因为在上帝里的恩典，而变成了甜蜜的生活。

\subsection{十二股水泉和七十棵树=十二+七十门徒}
出15:27“他们到了以琳，在那里有十二股水泉，七十棵棕树，他们就在那里的水边安营。”耶稣在公开的传道过程中，有两件事情也提到了十二和七十。路加福音第九章首先描写“耶稣叫齐了十二个门徒，给他们能力、权柄，制服一切的鬼，医治各样的病，又差遣他们去宣传上帝国的道，医治病人。”（路9：1-2）第十章耶稣继续地差派门徒去各方传播福音，到处治病。路10：1-3这事（就是差遣十二门徒）以后，主又设立了七十个人，差遣他们两个两个地在他前面，往自己所要到的各城、各地方去。就对他们说：“要收的庄稼多，作工的人少。所以你们当求庄稼的主，打发工人出去收他的庄稼。你们去吧！我差你们出去，如同羊羔进入狼群。”耶稣设立十二个门徒出去传福音之后，又设立了七十个人，同样地差遣他们出去传扬福音。这个十二和七十在某种意义上与以琳这里提到的十二股水泉和七十棵树相对应验了。

\subsection{吗哪}
第16章提到了以色列人在旷野赖以生存的一件重要的事——吗哪。他们来到汛的旷野，以色列人就发怨言，出16：3说：“巴不得我们早死在埃及地耶和华的手下，那时我们坐在肉锅旁边，吃得饱足；你们将我们领出来，在这旷野，是要叫这全会众都饿死啊！”这些人就抱怨说，在这么一个旷野，没吃没喝的，是要把我们全会众都饿死吗？倒不如我们坐在肉锅旁过把瘾就死还更痛快。这时，出16：4-5耶和华对摩西说：“我要将粮食从天降给你们。百姓可以出去，每天收每天的份，我好试验他们遵不遵我的法度。（如何试验的呢？就是）到第六天他们要把所收进来的预备好了，比每天所收的多一倍。”因为到第七天安息日，天上就不再降吗哪了。降吗哪这件事情象征我们属灵的生命，要靠耶稣——属天的吗哪来维系。

耶稣说：约6:48-51“我就是生命的粮。你们的祖宗在旷野吃过吗哪，还是死了。这是从天上降下来的粮，叫人吃了就不死，我是从天上降下来生命的粮；人若吃这粮，就必永远活着。我所要赐的粮，就是我的肉，为世人之生命所赐的。”这个肉怎么吃呢？犹太人就从字面上去谈这个问题。我们不要把圣经中的每一句话都过度的做灵意的解释，但是，圣经毕竟是圣灵感动的，是属灵的书籍，我们要带着属灵的眼光来读圣经。

约6:53-56耶稣说：“我实实在在地告诉你们: 你们若不吃人子的肉，不喝人子的血,就没有生命在你们里面。吃我肉，喝我血的人就有永生,在末日我要叫他复活。我的肉真是可吃的，我的血真是可喝的。吃我肉、喝我血的人常在我里面，我也常在他里面。”这是耶稣的一个吩咐。接着他说：约6:57-58“永活的父怎样差我来，我又因父活着；照样，吃我肉的人也要因我活着。这就是从天上降下来的粮。吃这粮的人，就永远活着，不像你们的祖宗吃过吗哪还是死了。”

\subsection{磐石}
出17：5-7 耶和华对摩西说：“你手里拿着你先前击打河水的杖，带领以色列的几个长老，从百姓面前走过去。我必在何烈的磐石那里站在你面前，你要击打磐石，从磐石里必有水流出来，使百姓可以喝。”摩西就在以色列的长老眼前这样行了。他给那地方起名叫玛撒，又叫米利巴，因以色列人争闹，又因他们试探耶和华说：“耶和华是在我们中间不是？”

这件事情之后，摩西击打磐石就有水流出来供以色列人来饮用。因上帝吩咐他说：你去拿着这个杖到何烈的磐石那里。那么这个磐石实际也就是对基督的一个预表。林前10：1-4“弟兄们，我不愿意你们不晓得，我们的祖宗从前都在云下，都从海中经过，都在云里、海里受洗归了摩西，并且都吃了一样的灵食，也都喝了一样的灵水；所喝的，是出于随着他们的灵磐石，那磐石就是基督。”所以这个出活水的磐石也就是基督的一个象征。


\subsection{祭司的国度，为圣洁的国民}
出19:2-6他们离了利非订，来到西奈的旷野，就在那里山下安营。摩西到上帝那里，耶和华从山上呼唤他说：“你要这样告诉雅各家，晓谕以色列人说：＇我向埃及人所行的事，你们都看见了；且看见我如鹰将你们背在翅膀上，带来归我。如今你们若实在听从我的话，遵守我的约，就要在万民中作属我的子民；因为全地都是我的，你们要归我作祭司的国度，为圣洁的国民。＇这些话你要告诉以色列人。” 在新约《彼得前书》2：9 “我们是圣洁的国度，君尊的祭司。”

\subsection{预备了身体，开通了耳朵}
我们在这看到，希伯来人如果是作奴仆的话，他服侍你六年，第七年他可以自由的走。如果在这中间他有结婚成家有生孩子，那么他第七年自由出去的时候他就可以自己离开，妻子和儿女要留下归给主人；可是如果他自愿的说我爱这个主人，我也舍不得和我的妻子儿女分开，我宁可不要自由，愿留下来继续地服侍，继续保持这个奴仆之身。那就要在他的耳朵上面给他穿一个洞，这就表明他永远地服侍这个主人了。

那么这个和耶稣有什么样的关系，为什么要订这样的一个条例呢？圣经上诗篇40：6-8祭物和礼物，你不喜悦，你已经开通我的耳朵。燔祭和赎罪祭非你所要。那时我说：“看哪，我来了！我的事在经卷上已经记载了。我的上帝啊，我乐意照你的旨意行，你的律法在我心里。”所以诗篇在40篇就说，上帝为这一位祂的圣子开通了他的耳朵，预备他来献上这个燔祭，他说他愿意照上帝的旨意行，上帝的律法在他的心里。同样的这段经文，在希伯来书10章也出现了，来10：5所以，基督到世上来的时候，就说：“上帝啊，祭物和礼物是你不愿意的，你曾给我预备了身体。”

我们注意到保罗说，上帝为耶稣“预备了身体”，而诗篇说“为他开通了耳朵”。那么我们在这儿看到，耶稣他本来是自由之身，他来到人间，他爱这个教会，不愿意自由离去，要永远的和他的妻子和儿子——他的教会在一起，愿意永远的服侍主人，那就在他的耳朵上开一个洞。所以圣经上说：你为我开通了耳朵。《出埃及记》21章，这个主人就带他到上帝那里，或者说审判官那里，又带他到门前，靠近这个门框的地方，用锥子穿他的耳朵。这件事情在诗篇中，也就是指上帝要为耶稣预备这样一个身体，他愿意永远地服侍主人，永远和他所爱的女人就是他的妻子——教会在一起。

\subsection{逃城}
我们同样的来看一下《出埃及记》21章后面的一些经文，出21：13-14“人若不是埋伏着杀人，乃是上帝交在他手中，我就设下一个地方，他可以往那里逃跑。人若任意用诡计杀了他的邻舍，就是逃到我的坛那里，也当捉去把他治死。”这里就说到有一个误杀的罪。如果有人误杀了人，就可以做一座逃城，让他暂时在这个逃城里边保命。圣经在后面记载，以色列实际建了六座逃城。


\subsection{牛触死了奴仆或者是婢女赔三十舍克勒}
出21：32，这里也有一节非常有趣的经文，一个条例。这个条例是这么说的：如果牛触了奴仆或者是婢女，就是你的牛把别人家里的家奴或者是婢女撞死，那怎么办呢？你就要将银子三十舍克勒给他们的主人，也要用石头把牛打死。（参出21：32）所以这里我们就注意到，这个奴仆和婢女他们的身价是多少呢？是三十舍克勒。这“三十舍克勒”马上就让我们想到耶稣被犹大出卖，他是多少钱出卖的呢？太26：15（犹大从耶稣面前退去，就去找祭司长找犹太公会说：）“我把他交给你们，你们愿意给我多少钱？”他们就给了他三十块钱。


\subsection{初熟之物/不可以用山羊羔母的奶来煮山羊羔}
我们来继续看一下出23：19，这节经文在申14：21，出34：26都有重复，那么重复的是什么呢？出23：19“地里首先初熟之物要送到耶和华你上帝的殿。不可用山羊羔母的奶煮山羊羔。”

这里两节都应验在耶稣的身上了。“地里首先初熟之物”，耶稣是初熟的果子，归给上帝。所以在这里可以说耶稣应验了这个“初熟之物”。下面这一句是我们在这要特别看到的，那就是“不可以用山羊羔母的奶来煮山羊羔”。请注意，如果你要煮山羊羔的话，说明这个山羊羔牠是已经被杀了。而山羊羔的奶，就是山羊羔的妈妈还是活着的，你要把山羊羔的妈妈和这个山羊羔分开，不要把牠们弄在一起，所以煮山羊羔之前要将山羊羔的妈妈——母亲分开。在耶稣被钉十字架的时候，耶稣的母亲和耶稣是分开的，没有在一起。耶稣的母亲并没有因耶稣被钉而受到法律的牵连。约19章记载耶稣的母亲在耶稣被钉时的情形。“站在耶稣十字架旁边的，有他母亲和他母亲的姐妹，并革罗罢的妻子马利亚和抹大拉的马利亚。耶稣见他母亲和他所爱的门徒站在旁边，就对他母亲说：“母亲，看，你的儿子！”又对那门徒说：“看，你的母亲！”从此那门徒就接她到自己家里去了。（约19：25-27）所以我们看到，这一节经文就预示着耶稣的母亲和这个山羊羔所预示的耶稣他们之间是分开的，山羊羔要被蒸煮，要死亡，可是母亲不会受到同样的律法的追究。

\subsection{差遣使者在你前面，在路上保护你，领你到我所预备的地方（十字架）}
出23：20看哪，我差遣使者在你前面，在路上保护你，领你到我所预备的地方去。

耶稣是上帝所差遣的，上帝要领他去哪里？上帝要领他去十字架这个地方，在路上要保护他，也差遣一个使者增加他的力量。路22：39-42耶稣出来，照常往橄榄山去，门徒也跟随。到了那地方，就对他们说：“你们要祷告，免得入了迷惑。”于是离开他们约有扔一块石头那么远，跪下祷告，说：“父啊，你若愿意，就把这杯撤去，然而，不要成就我的意思，只要成就你的意思。”在这个时候耶稣感觉非常痛苦，圣经上说，他是心里极其的伤痛，几乎要死。（可14：33-34）


\subsection{奉我名来的这一位可以赦罪}
出23:21他是奉我名来的，你们要在他面前谨慎，听从他的话，不可惹他，因为他必不赦免你们的过犯。这里说到这位要来的使者，他是奉主名来的，他有一个权柄就是给人赦罪，在可2章有一个很有名的故事，就是耶稣进了迦百农，很多人听见他在屋子里，就来了，结果把房子挤的是水泄不通。有一个瘫痪的人，他的四个朋友没有办法进去，就爬到房顶上，把瓦给揭开了，就把他从房顶上用绳子缒下来。耶稣看见他们的信心，就对瘫子说：“小子，你的罪赦了。”这对我们也是很有鼓励的一件事情。耶稣看见了他们的信心，他们是谁呢？是指这个瘫子呢？还是指着他的四个朋友，还是指着他们一起呢？

无论是哪一种情况，我们都可以看到，对那些病人，特别对那些无力帮助自己的人，我们可以用信心去帮助他们，上帝的美意，也可以在这样的人身上成就。所以耶稣说：“你的罪赦了。”因此这个被差遣的人是有赦罪的权柄的。这个故事里的这些文士坐在那里，心里议论说：“这个人为什么这样说呢？他说僭妄的话了，除了上帝以外，谁能赦罪呢？”耶稣知道他们心里这样议论，就说：“你们心里为什么这样议论呢？或对瘫子说＇你的罪赦了＇，或说＇起来，拿你的褥子行走＇，哪一样容易呢？但要叫你们知道，（知道什么？）人子在地上有赦罪的权柄。”（可2：1-10）这样就应验了出21章所说的，奉我名来的这一位，他是可以赦罪的。

\subsection{奉上帝的名而来的}
太21章耶稣最后荣进耶路撒冷，有很多人把棕榈枝摊在地上，把衣服也摊在地上，迎接耶稣荣耀进入耶路撒冷。太21：9前行后随的众人喊着说：“和散那归于大卫的子孙！奉主名来的，是应该称颂的！高高在上和散那！”在此我们看到，耶稣就是《出埃及记》中所说的“奉上帝的名而来的”。

\subsection{圣所}
\subsubsection{洗濯盆}
\subsubsection{铜祭坛}
\subsubsection{圣所，耶稣的一个预表——同在}
“又当为我造圣所，使我可以住在他们中间。制造帐幕和其中的一切器具，都要照我所指示你的样式。（出25:8-9）

\subsubsection{幔子}
在《出埃及记》26章看到殿里幔子的做法。请看出26:1 “你要用十幅幔子作帐幕。这些幔子要用捻的细麻和蓝色、紫色、朱红色线制造，并用巧匠的手工绣上基路伯。”这里请大家注意一下，这个幔子要用十幅，“十”让我们想起十条诫命，十个方面都是完美的。十幅组成这样一个幔子，幔子在圣经里代表着什么呢？来10:19-20  “弟兄们，我们既因耶稣的血得以坦然进入至圣所， 是藉着他给我们开了一条又新、又活的路从幔子经过，这幔子就是他的身体。”这个幔子就是预表着耶稣的身体。

\subsubsection{陈设饼}

\subsubsubsection{金灯台}


\subsubsection{金香炉}
\subsubsection{约柜}


\chapter{創世記}
\begin{table}[H]
\centering
\begin{tabular}{p{3cm}|p{8.cm}}\hline\hline
\multicolumn{1}{c|}{事件/人物}&\multicolumn{1}{c}{圣经章节}\\\hline
創造&1-2\\\cline{1-1}
犯罪&3-5\\\cline{1-1}
大洪水&6-9\\\cline{1-1}
迦南受诅和巴别塔&9-10\\\hline
\textcolor{red}{亚伯拉罕}&\textcolor{blue}{11-25（15章 30\%）}\\\hline
罗得&11:27-32/12:1-5/13/14:1-16/18:16-33/19（6章 12\%）\\\hline
以实马利&16/21/25（3章 6\%）\\\hline
\textcolor{red}{以撒}&\textcolor{blue}{21-35（15章 30\%）}\\\hline
以扫&25/32-33/36（4章 8\%）\\\hline
\textcolor{red}{雅各}&\textcolor{blue}{25-49（25章 50\%）}\\\hline
\textcolor{red}{约瑟}&\textcolor{blue}{30-50（21章 42\%）}\\\hline
犹大&37-38/42-44/49（6章 12\%）\\\hline
流便&\\\hline\hline
撒拉&11:27-32/12/16-18/20/21/23（8章 16\%）\\\hline
夏甲&16/21（2章 4\%）\\\hline
利百加&\\\hline
利雅&\\\hline
拉结&\\\hline
她玛&\\\hline
\end{tabular}
\end{table}


\section{神的創造 1-2}
\subsection{神創造天地和萬物 1:1-25}
\subsection{设立安息日 2:1-3}


\begin{table}[!hbp]
\centering
\begin{tabular}{p{1.1cm}|p{2.cm}| p{2.2cm}|p{2.2cm}|p{3cm}|p{3cm}}
%\begin{tabular}{|c|c|c|c|c|c|}
\hline
\hline
日期 & 創造 & 目的& 定義  & 創造方式&  结果 \\\hline
起初  & 神創造天地&& &神的靈運行在水面上&  地是空虛混沌，淵面黑暗 \\\hline
頭一日 & 光 & 把光暗分開& 光為晝,暗為夜&神說：要有光。就有了光&  神看光是好的  \\\hline
第二日&空氣 & 空氣以下空氣以上的水分開&  稱空氣為天& 神說：諸水之間要有空氣，將水分為上下& \\\hline
第三日&旱地 &天下的水聚在一處   &水的聚處為海 &神說：天下水要聚在一處，使旱地露出來& 神看著是好的\\
          &地發生青草&做走獸、飛鳥、爬在地上有生命物的食物　&　 &神說：地要發生青草&\\
          &結種子菜蔬結果子樹木　　　&賜給人作食物　&&和結種子的菜蔬，並結果子的樹木，\textcolor{blue}{各從其類}，果子都包著核&神看著是好的\\\hline
第四日&天上光體（兩個大光，眾星） &分晝夜、作記號、标记神的时间日期年．普照在地上&大的管晝，小的管夜&神說：天上要有光體，可以分晝夜，做記號，定節令、日子、年歲，並要發光在天空，普照在地上&神看著是好的\\\hline
第五日&水中有生命的魚和動物 && &神說：水要多多滋生\textcolor{blue}{各從其類}.神就\textcolor{red}{賜福}滋生繁多,充滿海 &神看著是好的。 \\
          & 天空的飛鳥 &  && 要有雀鳥飛在地上空中\textcolor{blue}{各從其類}.神就\textcolor{red}{賜福}  雀鳥多生&\\\hline
第六日& 地上活物（牲畜、昆蟲、野獸）& &&神說：地要生出活物來；牲畜,昆蟲,野獸\textcolor{blue}{各從其類}&神看著是好的\\
&照著神的形像、神的樣式造男造女 &要生養眾多，管理海裡的魚、空中的鳥、地上的牲畜、昆蟲 &&神就\textcolor{red}{賜福}給他們，又對他們說：要生養眾多，遍滿地面，治理這地，也要管理海裡的魚、空中的鳥和地上各樣行動的活物&神看著一切所造的都甚好\\\hline
第七日　&歇了他一切的工&安息&第七日定為聖日&神\textcolor{red}{賜福}給第七日，因為在這日神歇了他一切創造的工&天地萬物都造齊了，神造物的工已經完畢\\
\hline
\end{tabular}
\caption{神的創造}
\end{table}

\subsection{神造人 1:26-31，2:4-7，2:18-25}
\begin{itemize} 
\itemsep-.35em 
\item 创造计划
\begin{itemize} 
\itemsep-.35em 
\item 神說我們要照著我們的形象，按著我們的樣式造人, 神就照著自己的形象造男造女
\item 神就賜福給他們說：要生養眾多，遍滿地面。治理這地，也要管理海裡的魚、空中的鳥和地上各樣行動的活物
\item 神說：看哪，我將遍地上一切結種子的菜蔬和一切樹上所結有核的果子，全賜給你們做食物
\end{itemize}
\item 亞當受造\\
創造天地的來歷，在耶和華神造天地的日子，乃是這樣： 野地還沒有草木，田間的菜蔬還沒有長起來，因為耶和華神還沒有降雨在地上，也沒有人耕地。但有霧氣從地上騰，滋潤遍地。耶和華神用地上的塵土造人，將生氣吹在他鼻孔裡，他就成了有靈的活人，名叫亞當。
\item 夏娃受造
\begin{itemize} 
\itemsep-.35em 
\item 耶和華神說：「那人獨居不好，我要為他造一個配偶幫助他。」 
\item 耶和華神用土所造成的野地各樣走獸和空中各樣飛鳥都帶到那人面前看他叫什麼。那人怎樣叫各樣的活物，那就是牠的名字。那人便給一切牲畜和空中飛鳥野地走獸都起了名，只是那人沒有遇見配偶幫助他 
\item 耶和華神使他沉睡，他就睡了。於是取下他的一條肋骨，又把肉合起來。耶和華神就用那人身上所取的肋骨造成一個女人，領她到那人跟前。那人說：「這是我骨中的骨、肉中的肉！可以稱她為女人，因為她是從男人身上取出來的。」 
\end{itemize}
\item 建立婚姻——因此，人要離開父母，與妻子聯合，二人成為一體。當時夫妻二人赤身露體，並不羞恥
\end{itemize}

\subsection{神東方的伊甸園 2:8-17 }
\begin{itemize} 
\itemsep-.35em 
\item 人的职责——所造的人安置在那裡使他耕種和看守
\item 園内设施
\begin{itemize} 
\itemsep-.35em 
\item 有河從伊甸流出來，滋潤那園子，從那裡分為四道
\item 各樣的樹從地裡長出來，可以悅人的眼目，其上的果子好做食物。園子當中又有生命樹和分別善惡的樹
\end{itemize}
\item 神的吩咐
\begin{itemize} 
\itemsep-.35em 
\item 園中各樣樹上的果子，你可以隨意吃 
\item 只是分別善惡樹上的果子，你不可吃，因為你吃的日子必定死 
\end{itemize}
\end{itemize}

\begin{table}[H]
\centering
\begin{tabular}{|c|c|l|c|}
\hline
河順序   & 名稱 & 位置          & 出產                \\\hline
第一道河 & 比遜 & 環繞哈腓拉全地& 金子,珍珠和紅瑪瑙   \\\hline
第二道河 &基訓  & 環繞古實全地  &                     \\\hline
第三道河 &希底結-底格里斯&在亞述的東邊 &　						\\\hline
第四道河 &伯拉河-幼發拉底&賜給人作食物 &　						\\\hline
\end{tabular}
\caption{伊甸流出來的河}
\end{table}


\begin{table}[!hbp]
\centering
\begin{tabular}{|c|c|l|c| }
\hline
名稱  & 材料 & 創造方法   & 命名的作者           \\\hline
野地各樣走獸 &土所造成的 2:19&  &亞當   \\\hline
空中各樣飛鳥 &土所造成的  2:19&    &亞當       \\\hline
亞當 &地上的塵土&用塵土造人、將生氣吹在他鼻孔裡　&上帝	（受造後)					\\\hline
女人2:23/夏娃 &肋骨&神使他沉睡、取下他的一條肋骨、又把肉合起來 &亞當（受造後/犯罪後）						\\\hline
\end{tabular}
\caption{地上成員的來歷}
\end{table}




\section{始祖犯罪 3－5 }
\subsection{女人亚当犯罪 3:1-24 }
\subsubsection{被蛇誘惑 3:1-5 }


\subsubsection{違背主命 3:6-13 }



\subsubsection{犯罪代价 3:14-24 }
\begin{itemize} 
\itemsep-.35em 
\item 宣布审判
\begin{itemize} 
\itemsep-.35em 
\item 對蛇說：「你既做了這事，就必受咒詛，比一切的牲畜野獸更甚！你必用肚子行走，終身吃土。我又要叫你和女人彼此為仇，你的後裔和女人的後裔也彼此為仇，女人的後裔要傷你的頭，你要傷他的腳跟。」
\item 對女人說：「我必多多加增你懷胎的苦楚，你生產兒女必多受苦楚。你必戀慕你丈夫，你丈夫必管轄你。」
\item 對亞當說：「你既聽從妻子的話，吃了我所吩咐你不可吃的那樹上的果子，地必為你的緣故受咒詛，你必終身勞苦，才能從地裡得吃的。地必給你長出荊棘和蒺藜來，你也要吃田間的菜蔬。你必汗流滿面才得糊口，直到你歸了土，因為你是從土而出的。你本是塵土，仍要歸於塵土。」
\end{itemize}
\item 女人易名——亞當給他妻子起名叫夏娃，因為她是眾生之母
\item 神遮罪孽——耶和華神為亞當和他妻子用皮子做衣服給他們穿
\item 逐出伊甸 
\begin{itemize} 
\itemsep-.35em 
\item 耶和華神說：那人已經與我們相似，能知道善惡。現在恐怕他伸手又摘生命樹的果子吃，就永遠活著。便打發他出伊甸園去，耕種他所自出之土。於是把他趕出去了 
\item 在伊甸園的東邊安設基路伯和四面轉動發火焰的劍，要把守生命樹的道路
\end{itemize}
\end{itemize}
 
\begin{table}[H]
\centering
\begin{tabular}{|c|c|c|l|   }
\hline
\hline
樹名   & 說話的人 & 聽眾 &說的話                  \\\hline
分別善惡的樹&神& 亞當&園中各樣樹上的果子、你可以隨意喫．\\
            &  &   &只是分別善惡樹上的果子、你不可喫、因為你喫的日子必定死。\\
            &夏娃&蛇&園中樹上的果子我們可以喫。惟有園當中那棵樹上的果子、\\
            &    &&神曾說、你們不可喫、也不可摸、免得你們死。\\           
            &蛇  &女人&你們不一定死、因為　神知道、\\
            &    &&你們喫的日子眼睛就明亮了、你們便如　神能知道善惡。\\ \hline             
生命樹&神&神&那人已經與我們相似、能知道善惡．現在恐怕他\\
&&&伸手又摘生命樹的果子喫、就永遠活著．                  \\\hline
\end{tabular}
\caption{伊甸園的兩顆樹}
\end{table}

\iffalse
\begin{table}[!hbp]
\centering
\begin{tabular}{|c|c|c|c|    }
\hline
\hline
順序&調查對象   & 問的問題 & 回答                    \\\hline
1&那人       & 你在那裡。 & 我在園中聽見你的聲音、我就害怕、 \\
 &          &           &因為我赤身露體．我便藏了。\\\hline
 &          & 誰告訴你赤身露體呢、莫非你喫了 & 你所賜給我、與我同居的女人、\\
 &          & 我吩咐你不可喫的那樹上的果子麼& 他把那樹上的果子給我、我就喫了。\\\hline
2& 女人       &你作的是甚麼事呢。 &那蛇引誘我、我就喫了。\\\hline
\end{tabular}
\caption{神澄清事實的順序}
\end{table}


\begin{table}[!hbp]
\centering
\begin{tabular}{|c|l|     }
\hline
\hline
被審判的對象   & 審判                     \\\hline
蛇       & 你既作了這事、就必受咒詛、比一切的牲畜野獸更甚、
           你必用肚子行走、終身喫土。        \\\hline
           & 我又要叫你和女人彼此為仇、你的後裔和女人的後裔、也彼此為仇．\\
           & 也彼此為仇．女人的後裔要傷你的頭、你要傷他的腳跟。\\\hline
女人       &我必多多加增你懷胎的苦楚、你生產兒女必多受苦楚．\\
&你必戀慕你丈夫、你丈夫必管轄你。\\\hline
亞當 & 地必為你的緣故受咒詛．你必終身勞苦、纔能從地裡得喫的。\\
& 地必給你長出荊棘和蒺藜來、你也要喫田間的菜蔬。你必汗流滿面纔得糊口、\\
& 直到你歸了土、因為你是從土而出的．你本是塵土、仍要歸於塵土。\\\hline
\end{tabular}
\caption{神審判的順序}
\end{table}
\fi


\subsection{該隱犯罪——殺亞伯 4:1-24 }
\subsubsection{生亞伯該隱 4:1-5 }
\begin{table}[H]
\centering
\begin{tabular}{|l|l|l|l|l|l| }
\hline
姓名 & 顺序     &    职业 &    供物  &献祭的结果  &反应   \\\hline
該隱 &长子     &種地的  &地裡的出產&看不中\textcolor{red}{該隱}和他的供物  
&大大地發怒，變了臉色 \\\hline
亞伯&次子  &牧羊的 & 羊群中頭生的和羊的脂油&耶和華看中了\textcolor{red}{亞伯}和他的供物& \\\hline
\end{tabular}
\caption{亞伯及該隱}
\end{table}

\subsubsection{該隱殺弟被審判 4:6-15 }
\begin{table}[!hbp]
\centering
\begin{tabular}{|l|l|     }
\hline
\hline
上帝 &該隱                      \\\hline
你為甚麼發怒呢、你為甚麼變了臉色呢．    &該隱與他兄弟亞伯說話，二人正在田間   \\ 
你若行得好、豈不蒙悅納、你若行得不好、罪就伏在門前．  &該隱起來打他兄弟亞伯，把他殺了 \\
他必戀慕你、你卻要制伏他。           & \\\hline
你兄弟亞伯在哪裡？&我不知道。我豈是看守我兄弟的嗎？ \\\hline
你作了甚麼事呢、你兄弟的血、有聲音從地裡向我哀告。 &\\
地開了口、從你手裡接受你兄弟的血．現在你必從這地受咒詛。&\\ 
你種地、地不再給你效力．你必流離飄蕩在地上。&\\\hline
 & 我的刑罰太重、過於我所能當的。\\
& 你如今趕逐我離開這地、以致不見你面．\\
& 我必流離飄蕩在地上、凡遇見我的必殺我。\\\hline
凡殺該隱的必遭報七倍。就給該隱立一個記號，免得人遇見他就殺他&離開耶和華的面住在伊甸東挪得之地\\\hline
\end{tabular}
\caption{神審判該隱}
\end{table}

\subsubsection{該隱的家谱 4:16-24}
\begin{forest}
  for tree={
    child anchor=west,
    parent anchor=east,
    grow=east,
    draw,
    anchor=west,
    edge path={
      \noexpand\path[\forestoption{edge}]
        (.child anchor) -| +(-5pt,0) -- +(-5pt,0) |-
        (!u.parent anchor)\forestoption{edge label};
    },
  }
  [該隱-造以諾城
    [以諾 
    	[以拿 
    	 [米戶雅利[瑪土撒利[ 拉麥[亞大（妻）[雅八-住帳棚牲畜人祖師][猶八-彈琴吹簫人祖師]] [洗拉（妻）[土八該隱-銅匠鐵匠的祖師][拿瑪（女儿）]]]]]    	
      ]                                                   
    ]    
  ]
\end{forest}

\begin{table}[!hbp]
\centering
\begin{tabular}{|l|l|l|l|l| }
\hline
姓名 & 杀人前提    &    杀的人 &    发出咒诅的人  &  咒诅 \\\hline
該隱 &献祭不蒙悦纳     &兄弟亚伯 &耶和華 &  若殺該隱遭報七倍\\\hline
拉麥& 人傷我  &壯年人 & 拉麥&殺拉麥必遭報七十七倍 \\
拉麥& 人損我  &少年人 & 拉麥&殺拉麥必遭報七十七倍 \\\hline
\end{tabular}
\caption{拉麥對他兩個妻子說的咒诅和該隱的咒诅对比}
\end{table}


\subsection{亞當後裔的家譜 4:25-5:32}
\label{sec:YaDangJiaPu}
\subsubsection{亞當後裔的家譜 1 4:25-5:32}
\begin{tikzpicture}
\pgfplotstableread{ % Read the data into a table macro
Label   First   Second  Third
\textcolor{red}{}   1656   5 0
\textcolor{red}{}   1656   5 0
\textcolor{red}{}   1656   5 0
亞當      0     930     0. 
\textcolor{red}{}   1656   5 0
\textcolor{red}{}   1656   5 0
\textcolor{red}{}   1656   5 0
塞特      130   912     0. 
\textcolor{red}{}   1656   5 0
\textcolor{red}{}   1656   5 0
\textcolor{red}{}   1656   5 0
以挪士（開始求告神的名）   235   905     0. 
\textcolor{red}{}   1656   5 0
\textcolor{red}{}   1656   5 0
\textcolor{red}{}   1656   5 0
該南      325   910     0. 
\textcolor{red}{}   1656   5 0
\textcolor{red}{}   1656   5 0
\textcolor{red}{}   1656   5 0
瑪勒列    395   895     0
\textcolor{red}{}   1656   5 0
\textcolor{red}{}   1656   5 0
\textcolor{red}{}   1656   5 0
雅列      460   962     0
\textcolor{red}{}   1656   5 0
\textcolor{red}{}   1656   5 0
\textcolor{red}{}   1656   5 0
\textcolor{red}{\textbf{以諾}}      622   365     0
\textcolor{red}{}   1656   5 0
\textcolor{red}{}   1656   5 0
\textcolor{red}{}   1656   5 0
瑪土撒拉  687   969     0
\textcolor{red}{}   1656   5 0
\textcolor{red}{}   1656   5 0
\textcolor{red}{}   1656   5 0
拉麥      874   777    0
\textcolor{red}{}   1656   5 0
\textcolor{red}{}   1656   5 0
\textcolor{red}{}   1656   5 0
挪亞     1056   950   0
\textcolor{red}{}   1656   5 0
\textcolor{red}{}   1656   5 0
\textcolor{red}{}   1656   5 0
閃      1556   600   0
\textcolor{red}{}   1656   5 0
\textcolor{red}{}   1656   5 0
\textcolor{red}{}   1656   5 0
亞法撒  1656   438   0
\textcolor{red}{}   1656   5 0
\textcolor{red}{}   1656   5 0
\textcolor{red}{}   1656   5 0
沙拉  1691   433   0
\textcolor{red}{}   1656   5 0
\textcolor{red}{}   1656   5 0
\textcolor{red}{}   1656   5 0
希伯  1721   464   0
\textcolor{red}{}   1656   5 0
\textcolor{red}{}   1656   5 0
\textcolor{red}{}   1656   5 0
法勒  1755   239   0
\textcolor{red}{}   1656   5 0
\textcolor{red}{}   1656   5 0
\textcolor{red}{}   1656   5 0
拉吳  1785   239   0
\textcolor{red}{}   1656   5 0
\textcolor{red}{}   1656   5 0
\textcolor{red}{}   1656   5 0
西鹿  1817   230   0
\textcolor{red}{}   1656   5 0
\textcolor{red}{}   1656   5 0
\textcolor{red}{}   1656   5 0
拿鶴  1847   148   0
\textcolor{red}{}   1656   5 0
\textcolor{red}{}   1656   5 0
\textcolor{red}{}   1656   5 0
他拉  1876   205   0
\textcolor{red}{}   1656   5 0
\textcolor{red}{}   1656   5 0
\textcolor{red}{}   1656   5 0
\textcolor{blue}{\textbf{亞伯蘭}}  1946   175   0
\textcolor{red}{}   1656   5 0
\textcolor{red}{}   1656   5 0
\textcolor{red}{}   1656   5 0
\textcolor{blue}{\textbf{以撒}}  2046   180   0
\textcolor{red}{}   1656   5 0
\textcolor{red}{}   1656   5 0
\textcolor{red}{}   1656   5 0
\textcolor{blue}{\textbf{雅各}}  2106   147   0
\textcolor{red}{}   1656   5 0
\textcolor{red}{}   1656   5 0
\textcolor{red}{}   1656   5 0
}\datatable

\begin{axis}[
    xbar stacked,   % Stacked horizontal bars
    bar width=10pt,
    height=13. cm,
    width=15cm,
    xmin=0,         % Start x axis at 0 
    ytick=data,     % Use as many tick labels as y coordinates
    yticklabels from table={\datatable}{Label},  % Get the labels from the Label 
]
\addplot [fill=white,
        %draw=none,
        bar shift=1200pt,%remove the white bar part
                mark options={
            gray,
            thick
        },
        mark=waterfall bridge 1
	] 
	table [x=First, y expr=\coordindex] {\datatable};    % Plot the "First" column against the data index
\addplot [fill=green!70!blue]table [x=Second, y expr=\coordindex] {\datatable};
%\draw[,line width=1pt](160.5,400) -- (160.5,880);
\draw[,line width=1pt](200.5,400) -- (200.5,880);
\draw (160,980) node[below=3pt] {大洪水} node[above=3pt] {};
\draw (200,940) node[below=3pt] {挪亚死} node[above=3pt] { };
\draw[,line width=1pt](215.6,415) -- (215.6,920);
\draw (215,980) node[below=3pt] {闪死} node[above=3pt] { };
\end{axis}
\end{tikzpicture}


\subsubsection{亞當後裔的家譜 2(detailed) 4:25-5:32}
\begin{tikzpicture}
\pgfplotstableread{ % Read the data into a table macro
Label   First   Second  Third
\textcolor{red}{}   1656   5 0
\textcolor{red}{}   1656   5 0
\textcolor{red}{}   1656   5 0
亞當      0     930     0. 
\textcolor{red}{}   1656   5 0
\textcolor{red}{}   1656   5 0
\textcolor{red}{}   1656   5 0
塞特      130   912     0. 
\textcolor{red}{}   1656   5 0
\textcolor{red}{}   1656   5 0
\textcolor{red}{}   1656   5 0
以挪士   235   905     0. 
\textcolor{red}{}   1656   5 0
\textcolor{red}{}   1656   5 0
\textcolor{red}{}   1656   5 0
該南      325   910     0. 
\textcolor{red}{}   1656   5 0
\textcolor{red}{}   1656   5 0
\textcolor{red}{}   1656   5 0
瑪勒列    395   895     0
\textcolor{red}{}   1656   5 0
\textcolor{red}{}   1656   5 0
\textcolor{red}{}   1656   5 0
雅列      460   962     0
\textcolor{red}{}   1656   5 0
\textcolor{red}{}   1656   5 0
\textcolor{red}{}   1656   5 0
\textcolor{red}{\textbf{以諾}}      622   365     0
\textcolor{red}{}   1656   5 0
\textcolor{red}{}   1656   5 0
\textcolor{red}{}   1656   5 0
瑪土撒拉  687   969     0
\textcolor{red}{}   1656   5 0
\textcolor{red}{}   1656   5 0
\textcolor{red}{}   1656   5 0
拉麥      874   777    0
\textcolor{red}{}   1656   5 0
\textcolor{red}{}   1656   5 0
\textcolor{red}{}   1656   5 0
挪亞     1056   950   0
\textcolor{red}{}   1656   5 0
\textcolor{red}{}   1656   5 0
\textcolor{red}{}   1656   5 0
閃      1556   600   0
\textcolor{red}{}   1656   5 0
\textcolor{red}{}   1656   5 0
\textcolor{red}{}   1656   5 0
亞法撒  1656   438   0
\textcolor{red}{}   1656   5 0
\textcolor{red}{}   1656   5 0
\textcolor{red}{}   1656   5 0
沙拉  1691   433   0
\textcolor{red}{}   1656   5 0
\textcolor{red}{}   1656   5 0
\textcolor{red}{}   1656   5 0
希伯  1721   464   0
\textcolor{red}{}   1656   5 0
\textcolor{red}{}   1656   5 0
\textcolor{red}{}   1656   5 0
法勒  1755   239   0
\textcolor{red}{}   1656   5 0
\textcolor{red}{}   1656   5 0
\textcolor{red}{}   1656   5 0
拉吳  1785   239   0
\textcolor{red}{}   1656   5 0
\textcolor{red}{}   1656   5 0
\textcolor{red}{}   1656   5 0
西鹿  1817   230   0
\textcolor{red}{}   1656   5 0
\textcolor{red}{}   1656   5 0
\textcolor{red}{}   1656   5 0
拿鶴  1847   148   0
\textcolor{red}{}   1656   5 0
\textcolor{red}{}   1656   5 0
\textcolor{red}{}   1656   5 0
他拉  1876   205   0
\textcolor{red}{}   1656   5 0
\textcolor{red}{}   1656   5 0
\textcolor{red}{}   1656   5 0
\textcolor{blue}{\textbf{亞伯蘭}}  1946   175   0
\textcolor{red}{}   1656   5 0
\textcolor{red}{}   1656   5 0
\textcolor{red}{}   1656   5 0
\textcolor{blue}{\textbf{以撒}}  2046   180   0
\textcolor{red}{}   1656   5 0
\textcolor{red}{}   1656   5 0
\textcolor{red}{}   1656   5 0
\textcolor{blue}{\textbf{雅各}}  2106   147   0
\textcolor{red}{}   1656   5 0
\textcolor{red}{}   1656   5 0
\textcolor{red}{}   1656   5 0
}\datatable

\begin{axis}[
    xbar stacked,   % Stacked horizontal bars
    bar width=10pt,
    height=13. cm,
    width=17cm,
    xmin=0,         % Start x axis at 0
    ytick=data,     % Use as many tick labels as y coordinates
    yticklabels from table={\datatable}{Label},  % Get the labels from the Label column of the \datatable
]
\addplot [fill=white,
        %draw=none,
        bar shift=1200pt,%remove the white bar part
                mark options={
            gray,
            thick
        },
        mark=waterfall bridge 1
	] 
	table [x=First, y expr=\coordindex] {\datatable};    % Plot the "First" column against the data index
\addplot [fill=green!70!blue]table [x=Second, y expr=\coordindex] {\datatable};
\draw (106,70) node[below=3pt] {活930岁} node[above=3pt] {};
\draw (117,105) node[below=3pt] {活912岁} node[above=3pt] {};
\draw (148,142) node[below=3pt] {活905岁,人开始求告神名} node[above=3pt] {};
\draw (136,179) node[below=3pt] {活910岁} node[above=3pt] {};
\draw (141,222) node[below=3pt] {活895岁} node[above=3pt] {};
\draw (155,262) node[below=3pt] {活962岁} node[above=3pt] {};
\draw (145,304) node[below=3pt] {\textcolor{red}{與神同行300年，365岁神將他取去}} node[above=3pt] {};
\draw (205,342) node[below=3pt] {活969岁,大洪水同年死1656} node[above=3pt] {};
\draw (205,382) node[below=3pt] {活777岁,洪水5年前死1651} node[above=3pt] {};

\draw (214,423) node[below=3pt] {活950岁} node[above=3pt] {};
\draw (228,463) node[below=3pt] {活600岁} node[above=3pt] {};
\draw (228,502) node[below=3pt] {活438岁} node[above=3pt] {};
\draw (228,542) node[below=3pt] {活433岁} node[above=3pt] {};
\draw (231,582) node[below=3pt] {活464岁} node[above=3pt] {};
\draw (228,622) node[below=3pt] {活239岁} node[above=3pt] {};
\draw (228,662) node[below=3pt] {活239岁} node[above=3pt] {};

\draw (228,702) node[below=3pt] {活230岁} node[above=3pt] {};
\draw (228,742) node[below=3pt] {活148岁} node[above=3pt] {};
\draw (228,782) node[below=3pt] {活205岁} node[above=3pt] {};
\draw (228,822) node[below=3pt] {活175岁} node[above=3pt] {};
\draw (235,862) node[below=3pt] {活180岁} node[above=3pt] {};
\draw (237,902) node[below=3pt] {活147岁} node[above=3pt] {};

\draw (6, 105) node[below=3pt] {130} node[above=3pt] {};
\draw (16,141) node[below=3pt] {105} node[above=3pt] {};
\draw (28,179) node[below=3pt] {90} node[above=3pt] {};
\draw (35,219) node[below=3pt] {70} node[above=3pt] {};
\draw (42,257) node[below=3pt] {65} node[above=3pt] {};
\draw (55,300) node[below=3pt] {162} node[above=3pt] {};
\draw (65,338) node[below=3pt] {65} node[above=3pt] {};
\draw (82,379) node[below=3pt] {187} node[above=3pt] {};
\draw (99,421) node[below=3pt] {182} node[above=3pt] {};
\draw (151,468) node[below=3pt] {502} node[above=3pt] {};
\draw (160,507) node[below=3pt] {98} node[above=3pt] {};
\draw (80,755) node[below=3pt] {130+105+90+70+65+162+65+187+182+502+98=1656年} node[above=3pt] {};
\draw[>=triangle 45, <->] (0,700) -- (166,700);

\draw (141,900) node[below=3pt] {大洪水毁灭世界} node[above=3pt] {};
\draw (143,945) node[below=3pt] {亚当被造1656年} node[above=3pt] {};

\draw (193,953) node[below=3pt]{500年} node[above=3pt] {};
\draw[>=triangle 45, <->] (166,899) -- (216,899);

\draw[,line width=1pt](200.5,400) -- (200.5,815);%挪亚死
\draw (195,870) node[below=3pt] {挪亚死} node[above=3pt] { };

\draw[,line width=1pt](215.6,415) -- (215.6,920);%闪死
\draw (218,980) node[below=3pt] {闪死} node[above=3pt] { };
\draw[>=triangle 45, <-] (0,-88) -- (15,-45);
\draw (35,8) node[below=3pt] {亚当被造为0年} node[above=3pt] { };
\end{axis}
\end{tikzpicture}

\subsection{亞當和夏娃}
\subsubsection{神應許亞當和夏娃}
1:28神就賜福給他們，又對他們說：「要生養眾多，遍滿地面，治理這地，也要管理海裡的魚、空中的鳥和地上各樣行動的活物。」 

1:29 神說：「看哪，我將遍地上一切結種子的菜蔬和一切樹上所結有核的果子，全賜給你們做食物。 30 至於地上的走獸和空中的飛鳥，並各樣爬在地上有生命的物，我將青草賜給牠們做食物。」

2:16耶和華神吩咐他說：「園中各樣樹上的果子，你可以隨意吃； 17 只是分別善惡樹上的果子，你不可吃，因為你吃的日子必定死。」
2:18耶和華神說：「那人獨居不好，我要為他造一個配偶幫助他。」

3:16對女人說：「我必多多加增你懷胎的苦楚，你生產兒女必多受苦楚。你必戀慕你丈夫，你丈夫必管轄你。」 17 又對亞當說：「你既聽從妻子的話，吃了我所吩咐你不可吃的那樹上的果子，地必為你的緣故受咒詛，你必終身勞苦，才能從地裡得吃的。 18 地必給你長出荊棘和蒺藜來，你也要吃田間的菜蔬。 19 你必汗流滿面才得糊口，直到你歸了土，因為你是從土而出的。你本是塵土，仍要歸於塵土。」


\subsubsection{亞當和夏娃}
2:23 那人說：「這是我骨中的骨、肉中的肉！可以稱她為女人，因為她是從男人身上取出來的。」
3:20 亞當給他妻子起名叫夏娃，因為她是眾生之母。

6:3 耶和華說：「人既屬乎血氣，我的靈就不永遠住在他裡面，然而他的日子還可到一百二十年。」

\subsubsection{新约评价亞當和夏娃}


\section{挪亞及大洪水 6:1－11:9}
\subsection{發大洪水的原因 6:1－12}
\noindent
\begin{itemize} 
\itemsep-.35em 
\item 人的罪恶

\item 神的审判

\end{itemize}

\subsection{挪亞造方舟 6:13－22}%7:16}
\noindent
\begin{itemize} 
\itemsep-.35em 
\item \textbf{方舟要求} 
\begin{itemize} 
\itemsep-.35em 
\item 歌斐木造一隻方舟 
\item 分一間一間的造 
\item 裡外抹上松香 
\item 要長三百肘、寬五十肘、高三十肘 
\item 方舟上邊要留透光處、高一肘 
\item 方舟的門要開在旁邊 
\item 方舟要分上、中、下三層
\end{itemize}
\item \textbf{方舟裡的人，動物}
\begin{itemize} 
\itemsep-.35em 
\item 挪亞和他三個兒子、閃、含、雅弗、並挪亞的妻子、和三個兒婦、都進入方舟
\item 拿各樣食物積蓄起來、好作你和他們的食物
\item 凡有血肉、有氣息的活物、都一對一對的到挪亞那裡、進入方舟
\begin{itemize} 
\itemsep-.35em 
\item 不潔淨的畜類各從其類、你要帶一公一母 
\item 凡潔淨的畜類各從其類、你要帶七公七母 
\item 飛鳥各從其類，七公七母 
\item 地上的昆蟲各從其類、每樣兩個 
\end{itemize}
\end{itemize}
\end{itemize}

\subsection{大洪水 7:1－8:19 }
\noindent
\begin{itemize} 
\itemsep-.35em 
\item \underline{洪水前七天}，耶和華通知挪亚全家都要進入方舟
\begin{itemize} 
\itemsep-.35em 
\item 耶和華說：因為再過七天，我要降雨在地上四十晝夜，把我所造的各種活物都從地上除滅
\item 挪亞就同他的妻和兒子、兒婦，都進入方舟，躲避洪水
\item 潔淨的畜類和不潔淨的畜類，飛鳥並地上一切的昆蟲，都是一對一對有公有母，到挪亞那裡進入方舟
\item 挪亞就遵著耶和華所吩咐的行了。耶和華就把他關在方舟裡頭
\end{itemize}
\item 過了那七天洪水氾濫在地上。當挪亞六百歲\underline{二月十七日}那一天，大淵的泉源都裂開了，天上的窗戶也敞開了
\item \underline{四十晝夜}大雨降在地上 
\item 洪水氾濫在地上\underline{四十}天
\begin{itemize} 
\itemsep-.35em 
\item 水往上漲，把方舟從地上漂起。水勢浩大，在地上大大地往上漲，方舟在水面上漂來漂去
\item 水勢在地上極其浩大，天下的高山都淹沒了，水勢比山高過十五肘，山嶺都淹沒了
\item 凡在地上有血肉的動物，就是飛鳥、牲畜、走獸和爬在地上的昆蟲，以及所有的人都死了
\end{itemize}
\item 水勢浩大、在地上共\underline{一百五十天 }水就漸消 
\begin{itemize} 
\itemsep-.35em 
\item 神記念挪亞和挪亞方舟裡的一切走獸牲畜。神叫風吹地，水勢漸落
\item 淵源和天上的窗戶都閉塞了，天上的大雨也止住了。水從地上漸退．過了\underline{一百五十天}
\end{itemize}
\item \underline{七月十七日}方舟停在亞拉臘山上 
\item 水又漸消、到\underline{十月初一日}、山頂都現出來了 
\item 過了\underline{四十天}
\begin{itemize} 
\itemsep-.35em 
\item 挪亞開了方舟的窗戶放出一隻烏鴉去。那烏鴉飛來飛去，直到地上的水都乾了
\item 他又放出一隻鴿子去，要看看水從地上退了沒有。但遍地上都是水，鴿子找不著落腳之地，就回到方舟挪亞那裡，挪亞伸手把鴿子接進方舟來
\end{itemize}
\item 他又等了\underline{七天}、再把鴿子從方舟放出去．晚上鴿子叼著一個新擰下來的橄欖葉子，挪亞就知道地上的水退了
\item 他又等了\underline{七天}、放出鴿子去、鴿子就不再回來了 
\item 到挪亞六百零一歲,\underline{正月初一日}、地上的水都乾了 ,挪亞撤去方舟的蓋觀看，便見地面上乾了
\item \underline{二月二十七日}、地就都乾了。神吩咐挪亞出方舟 
\begin{itemize} 
\itemsep-.35em 
\item 你和你的妻子、兒子、兒婦都可以出方舟。在你那裡凡有血肉的活物，就是飛鳥、牲畜和一切爬在地上的昆蟲，都要帶出來，叫牠在地上多多滋生，大大興旺
\item 於是挪亞和他的妻子、兒子、兒婦都出來了。一切走獸、昆蟲、飛鳥和地上所有的動物，各從其類，也都出了方舟
\end{itemize} 
\end{itemize}

\subsection{挪亞築壇, 神與挪亞立約 8:20-9:29 }
\subsubsection{神與挪亞立約 8:20-9:17}
\begin{itemize} 
\itemsep-.35em 
\item 立約方式——築壇燔祭
\begin{itemize} 
\itemsep-.35em 
\item 挪亞為耶和華築了一座壇
\item 拿各類潔淨的牲畜、飛鳥、獻在壇上為燔祭
\end{itemize}
\item 立約對象
\begin{itemize} 
\itemsep-.35em 
\item 挪亞和挪亞的後裔立約 
\item 一切活物、就是飛鳥、牲畜、走獸．凡從方舟裡出來的活物 
\end{itemize}
\item 立約內容 
\begin{itemize} 
\itemsep-.35em 
\item 地還存留的時候、稼穡、寒暑、冬夏、晝夜、就永不停息了 
\item 挪亞和他的兒子們要生養眾多、遍滿了地 
\item 凡地上的走獸、和空中的飛鳥、都必驚恐、懼怕你們 
\item 連地上一切的昆蟲、並海裡一切的魚、都交付你們的手 
\item 凡活著的動物、都可以作你們的食物、這一切我都賜給你們如同菜蔬一樣 
\end{itemize}
\item 立約記號——彩虹
\begin{itemize} 
\itemsep-.35em 
\item 把虹放在雲彩中、這就可作我與地立約的記號了。我看見、就要記念所立的永約 
\item 凡有血肉的、不再被洪水滅絕、也不再有洪水毀壞地了 
\end{itemize}
\end{itemize}

\subsubsection{迦南受咒詛 9:18-29}
\begin{table}[H]
\centering
\begin{tabular}{c|l   }
\hline
\hline
兒子名稱   &     祝福或咒詛                 \\\hline
迦南    & 當受咒詛、必給他弟兄作奴僕的奴僕。 \\\hline
 閃          & 耶和華閃的　神、是應當稱頌的、願迦南作閃的奴僕。\\\hline
雅弗   &願　神使雅弗擴張、使他住在閃的帳棚裡．又願迦南作他的奴僕。\\\hline
\end{tabular}
\caption{迦南受咒詛}
\label{table:jianashouzu}
\end{table}

\subsection{挪亞的後裔 10:1－10:32 }
\noindent
\begin{forest}
  for tree={
    child anchor=west,
    parent anchor=east,
    grow=east,
    draw,
    anchor=west,
    edge path={
      \noexpand\path[\forestoption{edge}]
        (.child anchor) -| +(-5pt,0) -- +(-5pt,0) |-
        (!u.parent anchor)\forestoption{edge label};
    },
  }
  [挪亞[閃][含][雅弗]]    
  ]
\end{forest}

\begin{forest}
  for tree={
    child anchor=west,
    parent anchor=east,
    grow=east,
    draw,
    anchor=west,
    edge path={
      \noexpand\path[\forestoption{edge}]
        (.child anchor) -| +(-5pt,0) -- +(-5pt,0) |-
        (!u.parent anchor)\forestoption{edge label};
    },
  }
  [雅弗[歌篾[亞實基拿][利法][陀迦瑪]]
    	  [瑪各][瑪代]
    		[雅完[以利沙][他施][基提][多單]]
    		[土巴][米設][提拉] ]    
  ]
\end{forest}

\begin{forest}
  for tree={
    child anchor=west,
    parent anchor=east,
    grow=east,
    draw,
    anchor=west,
    edge path={
      \noexpand\path[\forestoption{edge}]
        (.child anchor) -| +(-5pt,0) -- +(-5pt,0) |-
        (!u.parent anchor)\forestoption{edge label};
    },
  }
  [含[古實[哈腓拉][撒弗他][西巴]
    			      [拉瑪[示巴][底但]]
              [寧錄（英勇的獵戶）[示拿地[巴別][以力][亞甲][甲尼]] 
    															        [亞述[尼尼微][利河伯][迦拉][利鮮]]]
    			[撒弗提迦]]
    			[麥西[路低人][亞拿米人][利哈比人][拿弗土希人]
    				  [帕斯魯細人][迦斯路希人][迦斐託人[非利士人]]]
    			[弗]
    			[迦南(西頓/基拉耳/所多瑪)[西頓][赫][耶布斯人][亞摩利人][革迦撒人][希未人]
    			     [亞基人][西尼人][亞瓦底人][洗瑪利人][哈馬人]]
	]   
  ]
\end{forest}

\begin{forest}
  for tree={
    child anchor=west,
    parent anchor=east,
    grow=east,
    draw,
    anchor=west,
    edge path={
      \noexpand\path[\forestoption{edge}]
        (.child anchor) -| +(-5pt,0) -- +(-5pt,0) |-
        (!u.parent anchor)\forestoption{edge label};
    },
  }
  [閃
    	[以攔][亞述][路德]
    	[亞法撒[沙拉[希伯[法勒（分的意思）][約坍[亞摩答][沙列][哈薩瑪非][耶拉]
    										  [哈多蘭][烏薩][德拉][俄巴路][亞比瑪利]
    										  [示巴][阿斐][哈腓拉][約巴]]]]]
    	[亞蘭[烏斯][戶勒][基帖][瑪施]]                                                   
    ] 
  ]
\end{forest}

\begin{figure} 
\centering
\includegraphics[width=6.in]{ZongJie/figs/AMoXiWuJing/Assyria}
\caption{亚述}
\label{fig:sub1}
\end{figure}

\begin{figure} 
\centering
\includegraphics[width=6.in]{ZongJie/figs/AMoXiWuJing/ChuangShiJi}
\caption{创世记}
\label{fig:sub1}
\end{figure}

\begin{figure} 
\centering
\includegraphics[width=6.in]{ZongJie/figs/AMoXiWuJing/ShanHanHouYi}
\caption{闪含后裔}
\label{fig:sub1}
\end{figure}

\subsubsection{神應許挪亞}
\subsubsection{挪亞應許子孙}
\subsubsection{新约评价挪亞}



\section{圣经人物}
\subsection{亞伯蘭——蒙召到迦南}
\subsubsection{亞伯拉罕家谱 11:10-32，22:20-24，25:1-6}
\label{sec:YLBLJiaPu}
\begin{forest}
  for tree={
    child anchor=west,
    parent anchor=east,
    grow=east,
    draw,
    anchor=west,
    edge path={
      \noexpand\path[\forestoption{edge}]
        (.child anchor) -| +(-5pt,0) -- +(-5pt,0) |-
        (!u.parent anchor)\forestoption{edge label};
    },
  }
  [他拉
    	[亞伯蘭[撒拉-妻，同父異母妹妹[以撒[雅各
    														[利亞-妻[流便1〔有兒子〕][西緬2〔聽見〕][利未3〔聯合〕][猶大4〔讚美〕][以薩迦9〔價值〕][西布倫10〔同住〕][底拿11（女兒）]]
    														[拉結（妻）[約瑟12〔增添〕][便雅憫13]]
    														[悉帕（大妾）[迦得7〔萬幸〕][亞設8〔有福〕]]
    														[辟拉（小妾）[但5〔伸冤〕][拿弗他利6〔相爭〕]]]
													[以掃]]]
    			 [夏甲（妾，使女）[以實瑪利]]
    			 [基土拉[心蘭][約珊[示巴][底但[亞書利族][利都是族][利烏米族]]][米但][米甸[以法][以弗][哈諾][亞比大][以勒大]][伊施巴][書亞]]]
    	[拿鶴[密迦（妻，哈蘭的女兒）[烏斯][布斯][基母利[亞蘭]][基薛][哈瑣][必達][益拉]
    										  [彼土利[利百加（以撒的妻子）][拉班
    										        [利亞（雅各的妻子）]
    										        [拉結（雅各的妻子）]]]]
			  [流瑪（妾）[提八][迦含][他轄][瑪迦]]]
    	[哈蘭（死在迦勒底的吾珥）[羅得[亞捫][摩押]] [密迦]                                                                                                  ]
    ] 
  ]
\end{forest}



\begin{sidewaystable} % <-- HERE

\begin{center}
\vspace{1cm}
\begin{minipage}{0.45\textwidth}
\tikzstyle{female} = [rectangle, draw, fill=red!20, rounded corners, minimum height=3em]
\tikzstyle{male} = [rectangle, draw, fill=blue!20, minimum height=3em]
\tikzstyle{neutral} = [rectangle, draw, fill=green!20, minimum height=3em]
\hspace{5cm}
\begin{tikzpicture}[level distance=4cm]
    \node[neutral] {\per{他拉}}
    [edge from parent fork right,grow=right]
        child {
            node[male] {\per{亞伯蘭(撒拉)}}             
              child {
                node[male] {\per{羅得}}
              } 
        }
        child {
            node[male] {\per{拿鶴(密迦)}}               
        };

\iffalse
    \node[neutral] {\per{他拉}}
    [edge from parent fork right,grow=right]
        child {     
            node[male] {\per{哈蘭}}     
            child {
                node[female] {\per{密迦}}
            }
            child {
                node[male] {\per{羅得}}
            }           
        }
        child {
            node[male] {\per{亞伯蘭}}
        }
        child {
            node[female] {\per{撒拉}}
        }
        child {
            node[male] {\per{拿鶴}}
        };
\fi
\end{tikzpicture}
\end{minipage}
\end{center}

\begin{center}
\includegraphics[scale=0.7, ]{ZongJie/figs/AMoXiWuJing/abrahammoving}
\end{center}


\vspace{1.5cm}
\hspace{4cm}
\subcaptionbox{}[.35\textwidth]{%
%\begin{minipage}{0.45\textwidth}
\begin{tikzpicture}[
  man/.style={rectangle,draw,fill=blue!20},
  woman/.style={rectangle,draw,fill=red!20,rounded corners=.8ex},
  grandchild/.style={grow=down,xshift=1em,anchor=west,
    edge from parent path={(\tikzparentnode.south) |- (\tikzchildnode.west)}},
  first/.style={level distance=6ex},
  second/.style={level distance=12ex},
  third/.style={level distance=18ex},
  level 1/.style={sibling distance=5em}]
    % Parents
    \coordinate
      child[grow=left] {node[man,anchor=east]{亞伯蘭}}
      child[grow=right] {node[woman,anchor=west]{撒拉}}
      child[grow=down,level distance=0ex]
    [edge from parent fork down]
    % Children and grandchildren
   % child{node[man] {Alfred}
    %  child[grandchild,first] {node[man]{Joe}}
    %  child[grandchild,second] {node[woman]{Heather}}
     % child[grandchild,third] {node[woman] {Barbara}}}
    %child{node[woman] {Berta}
     % child[grandchild,first] {node[man]{Howard}}};
    child {node[man] {羅得}};
   % child {node[woman]{Doris}
     % child[grandchild,first] {node[man]{Nick}}
    %  child[grandchild,second] {node[woman]{Liz}}};
\end{tikzpicture}}
%\end{minipage}

\vspace{-4cm}
\hspace{14cm}
%\begin{minipage}{0.45\textwidth}
\subcaptionbox{}[.5\textwidth]{%
\tikzstyle{female} = [rectangle, draw, fill=red!20, rounded corners, minimum height=3em]
\tikzstyle{male} = [rectangle, draw, fill=blue!20, minimum height=3em]
\tikzstyle{neutral} = [rectangle, draw, fill=green!20, minimum height=3em]
\hspace{5cm}
\begin{tikzpicture}[level distance=4cm]
 
    \node[neutral] {\per{他拉}}
    [edge from parent fork right,grow=right]
        child {
            node[male] {\per{亞伯蘭}}
        }
        child {
            node[female] {\per{撒拉}}
        }
        child {     
            node[male] {\per{哈蘭}}     
            child {
                node[male] {\per{羅得}}
            }  
            child {
                node[female] {\per{密迦}}
            }         
        }
        child {
            node[male] {\per{拿鶴}}
        };

\end{tikzpicture}}
%\end{minipage}
\end{sidewaystable} 


\subsubsection{亞伯蘭生平}

\begin{table}[H]
\begin{threeparttable}
%\begin{tabular}{@{}p{0.18\textwidth}*{4}{L{\dimexpr0.3\linewidth-8\tabcolsep\relax}}@{}}
\begin{tabular}{c c c c}
\toprule
 \multicolumn{1}{L{\dimexpr0.25\linewidth-1\tabcolsep\relax}}{\bfseries 事件} 
&\multicolumn{1}{L{\dimexpr0.2\linewidth-1\tabcolsep\relax}}{\bfseries  }   
&\multicolumn{1}{L{\dimexpr0.1\linewidth-1\tabcolsep\relax}}{\bfseries 圣经经文}     
&\multicolumn{1}{L{\dimexpr0.2\linewidth-1\tabcolsep\relax}}{\bfseries  备注}  \\\hline

\multicolumn{1}{L{\dimexpr0.25\linewidth-1\tabcolsep\relax}}{\bfseries 亞伯拉罕家谱}  
& \multicolumn{1}{L{\dimexpr0.25\linewidth-1\tabcolsep\relax}}{\bfseries  }
& \multicolumn{1}{L{\dimexpr0.1\linewidth-1\tabcolsep\relax}}{\bfseries 11:10-32}
& \multicolumn{1}{L{\dimexpr0.2\linewidth-1\tabcolsep\relax}}{\bfseries  }\\ 

\multicolumn{1}{L{\dimexpr0.25\linewidth-1\tabcolsep\relax}}{\bfseries 蒙召：吾珥-哈兰-迦南} 
& \multicolumn{1}{L{\dimexpr0.25\linewidth-1\tabcolsep\relax}}{\bfseries 離開本地、父家} 
& \multicolumn{1}{L{\dimexpr0.1\linewidth-1\tabcolsep\relax}}{\bfseries 12:1-5}
& \multicolumn{1}{L{\dimexpr0.2\linewidth-1\tabcolsep\relax}}{\bfseries 亞伯蘭75/撒莱65}\\ 

\multicolumn{1}{L{\dimexpr0.25\linewidth-1\tabcolsep\relax}}{\bfseries 在示劍/伯特利两次築壇} 
& \multicolumn{1}{L{\dimexpr0.25\linewidth-1\tabcolsep\relax}}{\bfseries 西邊是伯特利，東邊是艾}  
& \multicolumn{1}{L{\dimexpr0.1\linewidth-1\tabcolsep\relax}}{\bfseries 12:6-9}
& \multicolumn{1}{L{\dimexpr0.2\linewidth-1\tabcolsep\relax}}{\bfseries }\\

\multicolumn{1}{L{\dimexpr0.25\linewidth-1\tabcolsep\relax}}{\bfseries 因饑荒從南地下埃及} 
& \multicolumn{1}{L{\dimexpr0.255\linewidth-1\tabcolsep\relax}}{\bfseries 带埃及夏甲從南地去伯特利}   
& \multicolumn{1}{L{\dimexpr0.1\linewidth-1\tabcolsep\relax}}{\bfseries 12:10-20}
& \multicolumn{1}{L{\dimexpr0.2\linewidth-1\tabcolsep\relax}}{\bfseries }\\

\multicolumn{1}{L{\dimexpr0.25\linewidth-1\tabcolsep\relax}}{\bfseries 在伯特利與羅得分開} 
& \multicolumn{1}{L{\dimexpr0.2\linewidth-1\tabcolsep\relax}}{\bfseries 神向亞伯蘭重申應許} 
& \multicolumn{1}{L{\dimexpr0.1\linewidth-1\tabcolsep\relax}}{\bfseries 13:1-13}
& \multicolumn{1}{L{\dimexpr0.2\linewidth-1\tabcolsep\relax}}{\bfseries }\\

\multicolumn{1}{L{\dimexpr0.25\linewidth-1\tabcolsep\relax}}{\bfseries 在希伯倫築壇}  
& \multicolumn{1}{L{\dimexpr0.2\linewidth-1\tabcolsep\relax}}{\bfseries } 
& \multicolumn{1}{L{\dimexpr0.1\linewidth-1\tabcolsep\relax}}{\bfseries 13:14-18}
& \multicolumn{1}{L{\dimexpr0.2\linewidth-1\tabcolsep\relax}}{\bfseries }\\

\multicolumn{1}{L{\dimexpr0.25\linewidth-1\tabcolsep\relax}}{\bfseries 亞伯蘭救羅得} 
& \multicolumn{1}{L{\dimexpr0.3\linewidth-1\tabcolsep\relax}}{\bfseries 住在希伯崙幔利的橡樹}  
& \multicolumn{1}{L{\dimexpr0.1\linewidth-1\tabcolsep\relax}}{\bfseries 14:1-16}
& \multicolumn{1}{L{\dimexpr0.2\linewidth-1\tabcolsep\relax}}{\bfseries 见表~\ref{table:AbrahamSaveLot}}\\\hline

\multicolumn{1}{L{\dimexpr0.25\linewidth-1\tabcolsep\relax}}{\bfseries 献祭麥基洗德} 
& \multicolumn{1}{L{\dimexpr0.3\linewidth-1\tabcolsep\relax}}{\bfseries 沙微谷,王谷拒所多玛王财物}  
& \multicolumn{1}{L{\dimexpr0.1\linewidth-1\tabcolsep\relax}}{\bfseries 14:17-24}
& \multicolumn{1}{L{\dimexpr0.2\linewidth-1\tabcolsep\relax}}{\bfseries 见表~\ref{table:MaiJiXiDeSuoDuoMa} }\\

\multicolumn{1}{L{\dimexpr0.25\linewidth-1\tabcolsep\relax}}{\bfseries 神與亞伯蘭立約仪式} 
& \multicolumn{1}{L{\dimexpr0.2\linewidth-1\tabcolsep\relax}}{\bfseries 亞伯蘭彻底離開本族}  
& \multicolumn{1}{L{\dimexpr0.1\linewidth-1\tabcolsep\relax}}{\bfseries 15:1-21}
& \multicolumn{1}{L{\dimexpr0.2\linewidth-1\tabcolsep\relax}}{\bfseries 见表~\ref{table:YaBoLanYuGodLiYue}}\\

\multicolumn{1}{L{\dimexpr0.265\linewidth-1\tabcolsep\relax}}{\bfseries 撒萊將夏甲給亞伯蘭為妾}  
& \multicolumn{1}{L{\dimexpr0.2\linewidth-1\tabcolsep\relax}}{\bfseries } 
& \multicolumn{1}{L{\dimexpr0.1\linewidth-1\tabcolsep\relax}}{\bfseries 16:1-6}
& \multicolumn{1}{L{\dimexpr0.2\linewidth-1\tabcolsep\relax}}{\bfseries 亞伯拉罕85/撒拉75}\\

\multicolumn{1}{L{\dimexpr0.25\linewidth-1\tabcolsep\relax}}{\bfseries \textcolor{blue}{夏甲在書珥曠野遇見神}}  
& \multicolumn{1}{L{\dimexpr0.2\linewidth-1\tabcolsep\relax}}{\bfseries 加低斯，庇耳拉海萊} 
& \multicolumn{1}{L{\dimexpr0.1\linewidth-1\tabcolsep\relax}}{\bfseries 16:7-14}
& \multicolumn{1}{L{\dimexpr0.2\linewidth-1\tabcolsep\relax}}{\bfseries 见表~\ref{table:GodYuXiaJia}}\\

\multicolumn{1}{L{\dimexpr0.25\linewidth-1\tabcolsep\relax}}{\bfseries \textcolor{blue}{以實瑪利出生}}
& \multicolumn{1}{L{\dimexpr0.2\linewidth-1\tabcolsep\relax}}{\bfseries }   
& \multicolumn{1}{L{\dimexpr0.1\linewidth-1\tabcolsep\relax}}{\bfseries 16:15-16}
& \multicolumn{1}{L{\dimexpr0.2\linewidth-1\tabcolsep\relax}}{\bfseries 亞伯拉罕86/撒拉76}\\

\multicolumn{1}{L{\dimexpr0.25\linewidth-1\tabcolsep\relax}}{\bfseries 神給立約記號-割礼} 
& \multicolumn{1}{L{\dimexpr0.3\linewidth-1\tabcolsep\relax}}{\bfseries }  
& \multicolumn{1}{L{\dimexpr0.2\linewidth-1\tabcolsep\relax}}{\bfseries 17:1-14，22-27}
& \multicolumn{1}{L{\dimexpr0.2\linewidth-1\tabcolsep\relax}}{\bfseries 撒萊89,见表~\ref{table:GeLi}}\\

\multicolumn{1}{L{\dimexpr0.25\linewidth-1\tabcolsep\relax}}{\bfseries 神應許亞伯拉罕生以撒} 
& \multicolumn{1}{L{\dimexpr0.25\linewidth-1\tabcolsep\relax}}{\bfseries 亞伯蘭->亞伯拉罕99}  
& \multicolumn{1}{L{\dimexpr0.1\linewidth-1\tabcolsep\relax}}{\bfseries 17:15-21}
& \multicolumn{1}{L{\dimexpr0.2\linewidth-1\tabcolsep\relax}}{\bfseries 见表~\ref{table:GodPromisYBLH}}\\

\multicolumn{1}{L{\dimexpr0.25\linewidth-1\tabcolsep\relax}}{\bfseries 神應許撒拉生以撒} 
& \multicolumn{1}{L{\dimexpr0.2\linewidth-1\tabcolsep\relax}}{\bfseries 撒萊->撒拉89}  
& \multicolumn{1}{L{\dimexpr0.1\linewidth-1\tabcolsep\relax}}{\bfseries 18:1-15}
& \multicolumn{1}{L{\dimexpr0.2\linewidth-1\tabcolsep\relax}}{\bfseries 见表~\ref{table:GodPromiseIsaac}}\\

\multicolumn{1}{L{\dimexpr0.25\linewidth-1\tabcolsep\relax}}{\bfseries 亞伯拉罕為羅得代求}  
& \multicolumn{1}{L{\dimexpr0.2\linewidth-1\tabcolsep\relax}}{\bfseries 代求6次} 
& \multicolumn{1}{L{\dimexpr0.1\linewidth-1\tabcolsep\relax}}{\bfseries 18:16-33}
& \multicolumn{1}{L{\dimexpr0.2\linewidth-1\tabcolsep\relax}}{\bfseries 见表~\ref{table:YaBoLanPrayLuoDe}}\\

\multicolumn{1}{L{\dimexpr0.25\linewidth-1\tabcolsep\relax}}{\bfseries \textcolor{blue}{神滅所多瑪與蛾摩拉}} 
& \multicolumn{1}{L{\dimexpr0.25\linewidth-1\tabcolsep\relax}}{\bfseries 亞伯拉罕清早看其结局}  
& \multicolumn{1}{L{\dimexpr0.1\linewidth-1\tabcolsep\relax}}{\bfseries 19:1-29}
& \multicolumn{1}{L{\dimexpr0.2\linewidth-1\tabcolsep\relax}}{\bfseries }\\

\multicolumn{1}{L{\dimexpr0.25\linewidth-1\tabcolsep\relax}}{\bfseries \textcolor{blue}{摩押人亞捫人的起源}}  
& \multicolumn{1}{L{\dimexpr0.2\linewidth-1\tabcolsep\relax}}{\bfseries } 
& \multicolumn{1}{L{\dimexpr0.1\linewidth-1\tabcolsep\relax}}{\bfseries 19:30-38}
& \multicolumn{1}{L{\dimexpr0.2\linewidth-1\tabcolsep\relax}}{\bfseries }\\

\multicolumn{1}{L{\dimexpr0.25\linewidth-1\tabcolsep\relax}}{\bfseries 祝福亞比米勒-娶撒拉} 
& \multicolumn{1}{L{\dimexpr0.25\linewidth-1\tabcolsep\relax}}{\bfseries 加低斯,書珥中間的基拉耳}  
& \multicolumn{1}{L{\dimexpr0.1\linewidth-1\tabcolsep\relax}}{\bfseries 20:1-20}
& \multicolumn{1}{L{\dimexpr0.2\linewidth-1\tabcolsep\relax}}{\bfseries }\\

\multicolumn{1}{L{\dimexpr0.25\linewidth-1\tabcolsep\relax}}{\bfseries \textcolor{blue}{以撒出生}} 
& \multicolumn{1}{L{\dimexpr0.2\linewidth-1\tabcolsep\relax}}{\bfseries }  
& \multicolumn{1}{L{\dimexpr0.1\linewidth-1\tabcolsep\relax}}{\bfseries 21:1-7}
& \multicolumn{1}{L{\dimexpr0.2\linewidth-1\tabcolsep\relax}}{\bfseries }\\

\multicolumn{1}{L{\dimexpr0.25\linewidth-1\tabcolsep\relax}}{\bfseries 夏甲以實瑪利被逐}  
& \multicolumn{1}{L{\dimexpr0.2\linewidth-1\tabcolsep\relax}}{\bfseries 天使安慰夏甲} 
& \multicolumn{1}{L{\dimexpr0.1\linewidth-1\tabcolsep\relax}}{\bfseries 21:8-21}
& \multicolumn{1}{L{\dimexpr0.2\linewidth-1\tabcolsep\relax}}{\bfseries }\\\hline

\multicolumn{1}{L{\dimexpr0.25\linewidth-1\tabcolsep\relax}}{\bfseries 與亞比米勒立約} 
& \multicolumn{1}{L{\dimexpr0.2\linewidth-1\tabcolsep\relax}}{\bfseries 住別是巴}  
& \multicolumn{1}{L{\dimexpr0.1\linewidth-1\tabcolsep\relax}}{\bfseries 21:22-34}
& \multicolumn{1}{L{\dimexpr0.2\linewidth-1\tabcolsep\relax}}{\bfseries }\\ 

\multicolumn{1}{L{\dimexpr0.25\linewidth-1\tabcolsep\relax}}{\bfseries 亞伯拉罕獻以撒}  
& \multicolumn{1}{L{\dimexpr0.2\linewidth-1\tabcolsep\relax}}{\bfseries 住別是巴} 
& \multicolumn{1}{L{\dimexpr0.1\linewidth-1\tabcolsep\relax}}{\bfseries 22:1-19}
& \multicolumn{1}{L{\dimexpr0.2\linewidth-1\tabcolsep\relax}}{\bfseries }\\

\multicolumn{1}{L{\dimexpr0.25\linewidth-1\tabcolsep\relax}}{\bfseries 得知兄弟拿鶴家谱} 
& \multicolumn{1}{L{\dimexpr0.2\linewidth-1\tabcolsep\relax}}{\bfseries 利百加出生}  
& \multicolumn{1}{L{\dimexpr0.1\linewidth-1\tabcolsep\relax}}{\bfseries 22:20-24}
& \multicolumn{1}{L{\dimexpr0.2\linewidth-1\tabcolsep\relax}}{\bfseries }\\

\multicolumn{1}{L{\dimexpr0.25\linewidth-1\tabcolsep\relax}}{\bfseries 撒拉終葬在希伯崙}  
& \multicolumn{1}{L{\dimexpr0.25\linewidth-1\tabcolsep\relax}}{\bfseries 幔利就是希伯崙麥比拉洞} 
& \multicolumn{1}{L{\dimexpr0.1\linewidth-1\tabcolsep\relax}}{\bfseries 23:1-20}
& \multicolumn{1}{L{\dimexpr0.22\linewidth-1\tabcolsep\relax}}{\bfseries 撒拉127 见表~\ref{table:QiuDi}}\\

\multicolumn{1}{L{\dimexpr0.25\linewidth-1\tabcolsep\relax}}{\bfseries 以撒娶親} 
& \multicolumn{1}{L{\dimexpr0.25\linewidth-1\tabcolsep\relax}}{\bfseries 南地,庇耳拉海萊}  
& \multicolumn{1}{L{\dimexpr0.1\linewidth-1\tabcolsep\relax}}{\bfseries 24:1-67}
& \multicolumn{1}{L{\dimexpr0.22\linewidth-1\tabcolsep\relax}}{\bfseries 亞伯拉罕140/以撒40}\\

\multicolumn{1}{L{\dimexpr0.25\linewidth-1\tabcolsep\relax}}{\bfseries 亞伯拉罕家谱 3} 
& \multicolumn{1}{L{\dimexpr0.25\linewidth-1\tabcolsep\relax}}{\bfseries 妾基土拉的后裔}  
& \multicolumn{1}{L{\dimexpr0.1\linewidth-1\tabcolsep\relax}}{\bfseries 25:1-6}
& \multicolumn{1}{L{\dimexpr0.2\linewidth-1\tabcolsep\relax}}{\bfseries }\\

\multicolumn{1}{L{\dimexpr0.25\linewidth-1\tabcolsep\relax}}{\bfseries 亞伯拉罕離世}  
& \multicolumn{1}{L{\dimexpr0.25\linewidth-1\tabcolsep\relax}}{\bfseries 幔利就是希伯崙麥比拉洞} 
& \multicolumn{1}{L{\dimexpr0.1\linewidth-1\tabcolsep\relax}}{\bfseries 25:7-11}
& \multicolumn{1}{L{\dimexpr0.2\linewidth-1\tabcolsep\relax}}{\bfseries }\\

\multicolumn{1}{L{\dimexpr0.25\linewidth-1\tabcolsep\relax}}{\bfseries 以實瑪利後代} 
& \multicolumn{1}{L{\dimexpr0.2\linewidth-1\tabcolsep\relax}}{\bfseries 哈腓拉,書珥,亞述}  
& \multicolumn{1}{L{\dimexpr0.1\linewidth-1\tabcolsep\relax}}{\bfseries 25:12-18}
& \multicolumn{1}{L{\dimexpr0.2\linewidth-1\tabcolsep\relax}}{\bfseries }\\\bottomrule
\end{tabular}
\caption{亞伯拉罕生平}
\label{table:YaBoLanShengPing}
\end{threeparttable}
\end{table}

\begin{itemize} 
\itemsep-.35em 
\item 離開本地——出迦勒底的吾珥
\begin{itemize} 
\itemsep-.35em 
\item 創11:28 哈蘭死在他的本地迦勒底的吾珥，在他父親他拉之先
\item 創11:31 他拉帶著他兒子亞伯蘭，和他孫子哈蘭的兒子羅得，並他兒婦亞伯蘭的妻子撒萊，出了迦勒底的吾珥，要往迦南地去，他們走到哈蘭就住在那裡
\item 徒7:2 司提反說：「諸位父兄請聽！當日我們的祖宗亞伯拉罕在美索不達米亞，還未住哈蘭的時候，榮耀的神向他顯現，
\end{itemize}
\item 離開本族——出哈蘭到迦南
\begin{itemize} 
\itemsep-.35em 
\item 創12:4-5 亞伯蘭就照著耶和華的吩咐去了，羅得也和他同去。亞伯蘭出哈蘭的時候年75歲，亞伯蘭將他妻子撒萊和侄兒羅得，連他們在哈蘭所積蓄的財物人口，都帶往迦南地去
\item 創15:7 耶和華又對他說，我是耶和華，曾領你出了迦勒底的吾珥，為要將這地賜給你為業。
\item 尼9:7 你是耶和華神，曾揀選亞伯蘭，領他出迦勒底的吾珥，給他改名叫亞伯拉罕。
\item 徒7:4 他就離開迦勒底人之地，住在哈蘭。他父親死了以後，神使他從那裡搬到你們現在所住之地
\end{itemize}
\item 離開本家——与罗得分离
\begin{itemize} 
\itemsep-.35em 
\item 創12:6-7 亞伯蘭經過那地，到了示劍地方，摩利橡樹那裡，為向他顯現的耶和華築了一座壇
\item 創12:8 又遷到伯特利東邊的山支搭帳篷。西邊伯特利，東邊艾。他又為耶和華築壇求告耶和華的名
\item 創12:9 亞伯蘭又漸漸遷往南地去
\item 創12:10 那地遭遇饑荒。因饑荒甚大，亞伯蘭就下埃及去
\item 創13:1 從埃及上南地去。
\item 創13:3-4 亞伯蘭從南地漸漸往伯特利去，到了伯特利和艾的中間，就是他起先築壇的地方，他又在那裡求告耶和華的名
\item 創13:12 亞伯蘭在伯特利与罗得分离
\item 創13:18 亞伯蘭就搬了帳篷，來到希伯崙幔利的橡樹那裡居住，在那裡為耶和華築了一座壇
\item 創14:12 四王把亞伯蘭的侄兒羅得和羅得的財物擄掠去了，當時羅得正住在所多瑪。
\item 創14:15 亞伯蘭率領精練壯直追到但。又追到大馬士革左邊的何把，連他侄兒羅得和他的財物，以及婦女、人民也都奪回來
\item 創14:18 亞伯蘭殺敗回來的時候，所多瑪王出來，在沙微谷(王谷)迎接他,又有撒冷王麥基洗德帶著餅和酒出來迎接，亚伯兰献祭给撒冷王，拒绝所多玛王财物
\end{itemize}
\item 應許亞伯蘭肉身生子——拒绝所多玛王，选择撒冷王献祭
\begin{itemize} 
\itemsep-.35em 
\item 創15 神与亚伯兰举行立约仪式，應許亞伯蘭后裔生子
\item 創16 亞伯蘭85岁纳夏甲為妾，86岁生子以實瑪利
\end{itemize}
\item 應許亞伯拉罕妻子生以撒
\begin{itemize} 
\itemsep-.35em 
\item 創17 亞伯蘭99岁，神与亚伯蘭立约，改名，赐割礼，应许生以撒，给撒莱改名
\item 創18 亞伯拉罕在幔利橡樹接待三个天使，为所多玛代求
\item 創19：27 亞伯拉罕清早起來，到了他從前站在耶和華面前的地方，觀看所多瑪和蛾摩拉與平原
\item 創20 亞伯拉罕從那裡向南地遷去，寄居在加低斯和書珥中間的基拉耳，亞比米勒娶撒拉，亞伯拉罕为亞比米勒祝福
\item 創21 亞伯拉罕100岁得子以撒，以實瑪利被逐，亞比米勒與亞伯拉罕在別是巴立約，亞伯拉罕在非利士人的地寄居
\end{itemize}
\item 献上以撒
\begin{itemize} 
\itemsep-.35em 
\item 創22 神要亞伯拉罕在摩利亞山献以撒，并再次祝福亚伯拉罕，亞伯拉罕就住在別是巴，得知利百加出生
\item 創23 在基列亞巴（希伯崙）麥比拉，幔利前，赫人以弗崙的那塊田和其中的洞埋葬撒拉
\item 創23 差仆人回本族为以撒娶妻，以撒住在南地，庇耳拉海萊
\item 創24 亞伯拉罕繼娶基土拉生六子，并打發他們往東方去，175岁寿终 
\end{itemize}
\end{itemize}





\subsubsection{献祭麥基洗德，拒绝所多瑪王 14:17-24 见表~\ref{table:MaiJiXiDeSuoDuoMa}}
\begin{table}[H]
\begin{threeparttable}
\begin{tabular}{@{}p{0.1\textwidth}*{4}{L{\dimexpr0.45\linewidth-8\tabcolsep\relax}}@{}}
\toprule
 地點 &  王 & 亞伯拉罕              \\\hline

\multicolumn{1}{L{\dimexpr0.1\linewidth-2\tabcolsep\relax}}{\bfseries 撒冷王,沙微谷就是王谷} 
&\multicolumn{1}{L{\dimexpr0.45\linewidth-2\tabcolsep\relax}}{亞伯蘭殺敗基大老瑪和與他同盟的王回來的時候，\bfseries \textbf{撒冷王麥基洗德}帶著餅和酒出來迎接，他是至高神的祭司。他為亞伯蘭祝福，說：願天地的主、至高的神賜福於亞伯蘭！至高的神把敵人交在你手裡，是應當稱頌的！} 

&\multicolumn{1}{L{\dimexpr0.45\linewidth-2\tabcolsep\relax}}{\bfseries 亞伯蘭就把所得的拿出十分之一來，給麥基洗德} \\\hline

\multicolumn{1}{L{\dimexpr0.1\linewidth-2\tabcolsep\relax}}{\bfseries 索多瑪王,沙微谷就是王谷} 
&\multicolumn{1}{L{\dimexpr0.45\linewidth-2\tabcolsep\relax}}{亞伯蘭殺敗基大老瑪和與他同盟的王回來的時候，\bfseries \textbf{所多瑪王}出來、在沙微谷迎接他、沙微谷就是王谷。所多瑪王對亞伯蘭說：你把人口給我、財物你自己拿去罷。} 
&\multicolumn{1}{L{\dimexpr0.45\linewidth-2\tabcolsep\relax}}{\bfseries 亞伯蘭對所多瑪王說：我已經向天地的主、至高的　神耶和華起誓．凡是你的東西、就是一根線、一根鞋帶、我都不拿、免得你說、我使亞伯蘭富足．只有僕人所喫的、並與我同行的亞乃、以實各、幔利、所應得的分、可以任憑他們拿去。} \\\hline
\bottomrule
\end{tabular}
\caption{献祭麥基洗德，拒绝所多瑪王}
\label{table:MaiJiXiDeSuoDuoMa}
\end{threeparttable}
\end{table}

\subsubsection{新约对亚伯拉罕的评价}


\subsubsection{神在亚伯拉罕身上的工作过程 见表~\ref{table:GodWorkAbraham}}
\begin{table}[H]
\begin{threeparttable}
\begin{tabular}{@{}p{0.18\textwidth}*{4}{L{\dimexpr0.45\linewidth-8\tabcolsep\relax}}@{}}
\toprule
地點  &    神 & 亞伯拉罕              \\\hline
\multicolumn{1}{L{\dimexpr0.15\linewidth-2\tabcolsep\relax}}{\bfseries 哈蘭}  
& \multicolumn{1}{L{\dimexpr0.4\linewidth-2\tabcolsep\relax}}{\bfseries 你要離開本地、本族、父家、往我所要指示你的地去。我必叫你成為大國．我必賜福給你、叫你的名為大、你也要叫別人得福，為你祝福的、我必賜福與他、那咒詛你的、我必咒詛他、地上的萬族都要因你得福。} 
&\multicolumn{1}{L{\dimexpr0.35\linewidth-2\tabcolsep\relax}}{\bfseries 亞伯蘭將他妻子撒萊、和姪兒羅得、連他們在哈蘭所積蓄的財物、所得的人口、都帶往迦南地去} \\\hline

\multicolumn{1}{L{\dimexpr0.15\linewidth-2\tabcolsep\relax}}{\bfseries 示剑}  
& \multicolumn{1}{L{\dimexpr0.4\linewidth-2\tabcolsep\relax}}{\bfseries 耶和華向亞伯蘭說我要把這地賜給你的後裔}  & \multicolumn{1}{L{\dimexpr0.35\linewidth-2\tabcolsep\relax}}{\bfseries 亞伯蘭就在那裡為耶和華築了一座壇}\\\hline

\multicolumn{1}{L{\dimexpr0.15\linewidth-2\tabcolsep\relax}}{\bfseries 伯特利艾中間}  
&  \multicolumn{1}{L{\dimexpr0.4\linewidth-2\tabcolsep\relax}}{\bfseries }  
& \multicolumn{1}{L{\dimexpr0.35\linewidth-2\tabcolsep\relax}}{\bfseries 亞伯蘭就在那裡為耶和華築了一座壇}\\\hline

\multicolumn{1}{L{\dimexpr0.15\linewidth-2\tabcolsep\relax}}{\bfseries 南地}  
&  \multicolumn{1}{L{\dimexpr0.4\linewidth-2\tabcolsep\relax}}{\bfseries }  
& \multicolumn{1}{L{\dimexpr0.35\linewidth-2\tabcolsep\relax}}{\bfseries 遭遇饑荒}\\\hline

\multicolumn{1}{L{\dimexpr0.15\linewidth-2\tabcolsep\relax}}{\bfseries 埃及}  
&  \multicolumn{1}{L{\dimexpr0.4\linewidth-2\tabcolsep\relax}}{\bfseries 耶和華因亞伯蘭妻子撒萊的緣故、降大災與法老和他的全家。}  
& \multicolumn{1}{L{\dimexpr0.35\linewidth-2\tabcolsep\relax}}{\bfseries 法老吩咐人將亞伯蘭和他妻子、並他所有的都送走了}\\\hline

\multicolumn{1}{L{\dimexpr0.15\linewidth-2\tabcolsep\relax}}{\bfseries 南地}  
&  \multicolumn{1}{L{\dimexpr0.4\linewidth-2\tabcolsep\relax}}{\bfseries }  
& \multicolumn{1}{L{\dimexpr0.35\linewidth-2\tabcolsep\relax}}{\bfseries 亞伯蘭的金、銀、牲畜極多}\\\hline

\multicolumn{1}{L{\dimexpr0.15\linewidth-2\tabcolsep\relax}}{\bfseries 伯特利艾中間}  
&  \multicolumn{1}{L{\dimexpr0.4\linewidth-2\tabcolsep\relax}}{\bfseries  }  
& \multicolumn{1}{L{\dimexpr0.35\linewidth-2\tabcolsep\relax}}{\bfseries 迦南人、與比利洗人、在那地居住。亞伯蘭求告耶和華，和羅得彼此分離。}\\\hline

\multicolumn{1}{L{\dimexpr0.15\linewidth-2\tabcolsep\relax}}{\bfseries 迦南}  
&  \multicolumn{1}{L{\dimexpr0.4\linewidth-2\tabcolsep\relax}}{\bfseries  從你所在的地方、你舉目向東西南北觀看凡你所看見的一切地、我都要賜給你和你的後裔、直到永遠．我也要使你的後裔如同地上的塵沙那樣多、人若能數算地上的塵沙、纔能數算你的後裔。你起來、縱橫走遍這地、因為我必把這地賜給你。}  
& \multicolumn{1}{L{\dimexpr0.35\linewidth-2\tabcolsep\relax}}{\bfseries }\\\hline

\multicolumn{1}{L{\dimexpr0.15\linewidth-2\tabcolsep\relax}}{\bfseries 希伯崙幔利的橡樹那}  
&  \multicolumn{1}{L{\dimexpr0.4\linewidth-2\tabcolsep\relax}}{\bfseries }  
& \multicolumn{1}{L{\dimexpr0.35\linewidth-2\tabcolsep\relax}}{\bfseries 為耶和華築了一座壇。}\\\hline

\multicolumn{1}{L{\dimexpr0.15\linewidth-2\tabcolsep\relax}}{\bfseries 亞摩利人幔利的橡樹}  
&  \multicolumn{1}{L{\dimexpr0.4\linewidth-2\tabcolsep\relax}}{\bfseries }  
& \multicolumn{1}{L{\dimexpr0.35\linewidth-2\tabcolsep\relax}}{\bfseries 亞伯蘭率領他家裡生養的精練壯丁三百一十八人去救羅得。}\\\hline

\multicolumn{1}{L{\dimexpr0.15\linewidth-2\tabcolsep\relax}}{\bfseries 但}  
&  \multicolumn{1}{L{\dimexpr0.4\linewidth-2\tabcolsep\relax}}{\bfseries }  
& \multicolumn{1}{L{\dimexpr0.35\linewidth-2\tabcolsep\relax}}{\bfseries 夜間、自己同僕人分隊殺敗敵人。}\\\hline
\multicolumn{1}{L{\dimexpr0.15\linewidth-2\tabcolsep\relax}}{\bfseries 大馬色左邊的何把}  
&  \multicolumn{1}{L{\dimexpr0.4\linewidth-2\tabcolsep\relax}}{\bfseries }  
& \multicolumn{1}{L{\dimexpr0.35\linewidth-2\tabcolsep\relax}}{\bfseries 將被擄掠的一切財物、連他姪兒羅得和他的財物、以及婦女人民都奪回來。}\\\hline

\multicolumn{1}{L{\dimexpr0.15\linewidth-2\tabcolsep\relax}}{\bfseries 沙微谷就是王谷}  
&  \multicolumn{1}{L{\dimexpr0.4\linewidth-2\tabcolsep\relax}}{\bfseries 撒冷王麥基洗德、帶著餅和酒、出來迎接．他是至高　神的祭司。他為亞伯蘭祝福、說、願天地的主、至高的神、賜福與亞伯蘭．至高的神把敵人交在你手裡、是應當稱頌的。 }  
& \multicolumn{1}{L{\dimexpr0.35\linewidth-2\tabcolsep\relax}}{\bfseries 亞伯蘭就把所得的、拿出十分之一來、給麥基洗德。  }\\\hline

\multicolumn{1}{L{\dimexpr0.15\linewidth-2\tabcolsep\relax}}{\bfseries ？}  
&  \multicolumn{1}{L{\dimexpr0.4\linewidth-2\tabcolsep\relax}}{\bfseries 耶和華與亞伯蘭立約。 }  
& \multicolumn{1}{L{\dimexpr0.35\linewidth-2\tabcolsep\relax}}{\bfseries  }\\\hline

\multicolumn{1}{L{\dimexpr0.15\linewidth-2\tabcolsep\relax}}{\bfseries ？}  
&  \multicolumn{1}{L{\dimexpr0.4\linewidth-2\tabcolsep\relax}}{\bfseries 亞伯蘭年八十六歲生以實瑪利 }  
& \multicolumn{1}{L{\dimexpr0.35\linewidth-2\tabcolsep\relax}}{\bfseries }\\\hline

\multicolumn{1}{L{\dimexpr0.15\linewidth-2\tabcolsep\relax}}{\bfseries 幔利橡樹}  
&  \multicolumn{1}{L{\dimexpr0.4\linewidth-2\tabcolsep\relax}}{\bfseries 耶和華應許生以撒 }  
& \multicolumn{1}{L{\dimexpr0.35\linewidth-2\tabcolsep\relax}}{\bfseries 亞伯拉罕九十九歲受割禮 }\\\hline

\multicolumn{1}{L{\dimexpr0.15\linewidth-2\tabcolsep\relax}}{\bfseries 基拉耳}  
&  \multicolumn{1}{L{\dimexpr0.4\linewidth-2\tabcolsep\relax}}{\bfseries }  
& \multicolumn{1}{L{\dimexpr0.35\linewidth-2\tabcolsep\relax}}{\bfseries 亞伯拉罕為基拉耳王禱告}\\\hline

\multicolumn{1}{L{\dimexpr0.15\linewidth-2\tabcolsep\relax}}{\bfseries ？}  
&  \multicolumn{1}{L{\dimexpr0.4\linewidth-2\tabcolsep\relax}}{\bfseries 你不必為這童子和你的使女憂愁、凡撒拉對你說的話、你都該聽從．因為從以撒生的、纔要稱為你的後裔。至於使女的兒子、我也必使他的後裔成立一國、因為他是你所生的。}  
& \multicolumn{1}{L{\dimexpr0.35\linewidth-2\tabcolsep\relax}}{\bfseries 以撒斷奶的日子趕走以實瑪利}\\\hline

\bottomrule
\end{tabular}
\caption{神在亞伯拉罕身上的工作过程}
\label{table:GodWorkAbraham}
\end{threeparttable}
\end{table}


%\subsection{亞伯拉罕與亞比米勒 20, 21:22-34 }
\subsubsection{基拉耳:亞伯拉罕為亞比米勒祝福 20 见表~\ref{table:ZhuFuYBML}}
\begin{table}[H]
\begin{threeparttable}
\begin{tabular}{@{}p{0.1\textwidth}*{3}{L{\dimexpr0.12\linewidth-2\tabcolsep\relax}}@{}}
%\begin{tabular}{@{}p{0.18\textwidth}*{4}{L{\dimexpr0.45\linewidth-8\tabcolsep\relax}}@{}}
\toprule
   亞伯拉罕  & 亞比米勒   &    神     \\\hline
\multicolumn{1}{L{\dimexpr0.3\linewidth-2\tabcolsep\relax}}{\bfseries 亞伯拉罕從那裡向南地遷去，寄居在加低斯和書珥中間的基拉耳。亞伯拉罕稱他的妻撒拉為妹子} 
&\multicolumn{1}{L{\dimexpr0.45\linewidth-2\tabcolsep\relax}}{\bfseries 基拉耳王亞比米勒差人把撒拉取了去} 
&\multicolumn{1}{L{\dimexpr0.25\linewidth-2\tabcolsep\relax}}{\bfseries 但夜間神來在夢中對亞比米勒說：「你是個死人哪！因為你取了那女人來，她原是別人的妻子。」} \\

\multicolumn{1}{L{\dimexpr0.3\linewidth-2\tabcolsep\relax}}{\bfseries } 
&\multicolumn{1}{L{\dimexpr0.45\linewidth-2\tabcolsep\relax}}{\bfseries } 
&\multicolumn{1}{L{\dimexpr0.25\linewidth-2\tabcolsep\relax}}{\bfseries } \\

\multicolumn{1}{L{\dimexpr0.3\linewidth-2\tabcolsep\relax}}{\bfseries } 
&\multicolumn{1}{L{\dimexpr0.45\linewidth-2\tabcolsep\relax}}{\bfseries 亞比米勒卻還沒有親近撒拉，他說：「主啊，連有義的國，你也要毀滅嗎？那人豈不是自己對我說『她是我的妹子』嗎？就是女人也自己說『他是我的哥哥』。我做這事是心正手潔的。」} 
&\multicolumn{1}{L{\dimexpr0.25\linewidth-2\tabcolsep\relax}}{\bfseries 神在夢中對他說：「我知道你做這事是心中正直，我也攔阻了你，免得你得罪我。所以我不容你沾著她。現在你把這人的妻子歸還他，因為他是先知，他要為你禱告，使你存活。你若不歸還他，你當知道，你和你所有的人都必要死。」} \\

\multicolumn{1}{L{\dimexpr0.3\linewidth-2\tabcolsep\relax}}{\bfseries } 
&\multicolumn{1}{L{\dimexpr0.45\linewidth-2\tabcolsep\relax}}{\bfseries 亞比米勒清早起來，召了眾臣僕來，將這些事都說給他們聽，他們都甚懼怕。亞比米勒召了亞伯拉罕來，對他說：「你怎麼向我這樣行呢？我在什麼事上得罪了你，你竟使我和我國裡的人陷在大罪裡？你向我行不當行的事了。」亞比米勒又對亞伯拉罕說：「你見了什麼才做這事呢？」}
&\multicolumn{1}{L{\dimexpr0.25\linewidth-2\tabcolsep\relax}}{\bfseries  } \\

\multicolumn{1}{L{\dimexpr0.3\linewidth-2\tabcolsep\relax}}{\bfseries } 
&\multicolumn{1}{L{\dimexpr0.45\linewidth-2\tabcolsep\relax}}{\bfseries } 
&\multicolumn{1}{L{\dimexpr0.25\linewidth-2\tabcolsep\relax}}{\bfseries } \\


\multicolumn{1}{L{\dimexpr0.3\linewidth-2\tabcolsep\relax}}{\bfseries 亞伯拉罕說：「我以為這地方的人總不懼怕神，必為我妻子的緣故殺我。況且她也實在是我的妹子，她與我是同父異母，後來做了我的妻子。當神叫我離開父家，漂流在外的時候，我對她說：『我們無論走到什麼地方，你可以對人說「他是我的哥哥」，這就是你待我的恩典了。』」} 
&\multicolumn{1}{L{\dimexpr0.45\linewidth-2\tabcolsep\relax}}{\bfseries 亞比米勒把牛、羊、僕婢賜給亞伯拉罕，又把他的妻子撒拉歸還他。亞比米勒又說：「看哪，我的地都在你面前，你可以隨意居住。」又對撒拉說：「我給你哥哥一千銀子，作為你在閤家人面前遮羞(眼)的，你就在眾人面前沒有不是了。」} 
&\multicolumn{1}{L{\dimexpr0.25\linewidth-2\tabcolsep\relax}}{\bfseries } \\

\multicolumn{1}{L{\dimexpr0.3\linewidth-2\tabcolsep\relax}}{\bfseries } 
&\multicolumn{1}{L{\dimexpr0.45\linewidth-2\tabcolsep\relax}}{\bfseries } 
&\multicolumn{1}{L{\dimexpr0.25\linewidth-2\tabcolsep\relax}}{\bfseries } \\

\multicolumn{1}{L{\dimexpr0.3\linewidth-2\tabcolsep\relax}}{\bfseries 亞伯拉罕禱告神} 
&\multicolumn{1}{L{\dimexpr0.45\linewidth-2\tabcolsep\relax}}{\bfseries } 
&\multicolumn{1}{L{\dimexpr0.25\linewidth-2\tabcolsep\relax}}{\bfseries 神就醫好了亞比米勒和他的妻子，並他的眾女僕，她們便能生育。因耶和華為亞伯拉罕的妻子撒拉的緣故，已經使亞比米勒家中的婦人不能生育}\\
\bottomrule
\end{tabular}
\caption{基拉耳:為亞比米勒祝福 }
\label{table:ZhuFuYBML}
\end{threeparttable}
\end{table}





\subsubsection{別是巴:亞伯拉罕與亞比米勒立約 21:22-34 见表~\ref{table:YaBoLanLiYue}}
\begin{table}[H]
\begin{threeparttable}
\begin{tabular}{@{}p{0.4\textwidth}*{4}{L{\dimexpr0.45\linewidth-8\tabcolsep\relax}}@{}}
\toprule
  亞比米勒王 & 亞伯拉罕              \\\hline


%\multicolumn{1}{L{\dimexpr0.1\linewidth-2\tabcolsep\relax}}{\bfseries  別是巴} 
\multicolumn{1}{L{\dimexpr0.45\linewidth-2\tabcolsep\relax}}{\bfseries 當那時候，亞比米勒同他軍長非各對亞伯拉罕說：「凡你所行的事都有神的保佑。我願你如今在這裡指著神對我起誓，不要欺負我與我的兒子，並我的子孫。我怎樣厚待了你，你也要照樣厚待我與你所寄居這地的民。」} 
&\multicolumn{1}{L{\dimexpr0.45\linewidth-2\tabcolsep\relax}}{\bfseries 亞伯拉罕說：「我情願起誓。」} \\ 

%\multicolumn{1}{L{\dimexpr0.1\linewidth-2\tabcolsep\relax}}{\bfseries } 
\multicolumn{1}{L{\dimexpr0.45\linewidth-2\tabcolsep\relax}}{\bfseries } 
&\multicolumn{1}{L{\dimexpr0.45\linewidth-2\tabcolsep\relax}}{\bfseries } \\

%\multicolumn{1}{L{\dimexpr0.1\linewidth-2\tabcolsep\relax}}{\bfseries  } 
 \multicolumn{1}{L{\dimexpr0.45\linewidth-2\tabcolsep\relax}}{\bfseries 從前亞比米勒的僕人霸占了一口水井} 
&\multicolumn{1}{L{\dimexpr0.45\linewidth-2\tabcolsep\relax}}{\bfseries 亞伯拉罕為這事指責亞比米勒。} \\ 

%\multicolumn{1}{L{\dimexpr0.1\linewidth-2\tabcolsep\relax}}{\bfseries } 
 \multicolumn{1}{L{\dimexpr0.45\linewidth-2\tabcolsep\relax}}{\bfseries } 
&\multicolumn{1}{L{\dimexpr0.45\linewidth-2\tabcolsep\relax}}{\bfseries } \\


%\multicolumn{1}{L{\dimexpr0.1\linewidth-2\tabcolsep\relax}}{\bfseries  } 
 \multicolumn{1}{L{\dimexpr0.45\linewidth-2\tabcolsep\relax}}{\bfseries 亞比米勒說：「誰做這事，我不知道，你也沒有告訴我，今日我才聽見了。」} 
&\multicolumn{1}{L{\dimexpr0.45\linewidth-2\tabcolsep\relax}}{\bfseries 亞伯拉罕把羊和牛給了亞比米勒，二人就彼此立約。亞伯拉罕把七隻母羊羔另放在一處} \\ 

% \multicolumn{1}{L{\dimexpr0.1\linewidth-2\tabcolsep\relax}}{\bfseries } 
 \multicolumn{1}{L{\dimexpr0.45\linewidth-2\tabcolsep\relax}}{\bfseries } 
&\multicolumn{1}{L{\dimexpr0.45\linewidth-2\tabcolsep\relax}}{\bfseries } \\


%\multicolumn{1}{L{\dimexpr0.1\linewidth-2\tabcolsep\relax}}{\bfseries  } 
 \multicolumn{1}{L{\dimexpr0.45\linewidth-2\tabcolsep\relax}}{\bfseries 亞比米勒問亞伯拉罕說：「你把這七隻母羊羔另放在一處，是什麼意思呢？」} 
&\multicolumn{1}{L{\dimexpr0.45\linewidth-2\tabcolsep\relax}}{\bfseries 他說：「你要從我手裡受這七隻母羊羔，做我挖這口井的證據。」所以他給那地方起名叫別是巴(盟誓的井)，因為他們二人在那裡起了誓} \\ 

%\multicolumn{1}{L{\dimexpr0.1\linewidth-2\tabcolsep\relax}}{\bfseries } 
 \multicolumn{1}{L{\dimexpr0.45\linewidth-2\tabcolsep\relax}}{\bfseries } 
&\multicolumn{1}{L{\dimexpr0.45\linewidth-2\tabcolsep\relax}}{\bfseries } \\


%\multicolumn{1}{L{\dimexpr0.1\linewidth-2\tabcolsep\relax}}{\bfseries  } 
 \multicolumn{1}{L{\dimexpr0.45\linewidth-2\tabcolsep\relax}}{\bfseries 他們在別是巴立了約，亞比米勒就同他軍長非各起身回非利士地去了} 
&\multicolumn{1}{L{\dimexpr0.45\linewidth-2\tabcolsep\relax}}{\bfseries 亞伯拉罕在別是巴栽上一棵垂絲柳樹，又在那裡求告耶和華永生神的名。亞伯拉罕在非利士人的地寄居了多日} \\ 
\bottomrule
\end{tabular}
\caption{別是巴:與亞比米勒立約}
\label{table:YaBoLanLiYue}
\end{threeparttable}
\end{table}



\subsection{亞伯蘭侄子羅得——跟随亚伯蘭到迦南}
\subsubsection{羅得生平}
\begin{table}[H]
\centering
\begin{tabular}{p{14cm}|p{2.350cm}}\hline\hline
\multicolumn{1}{c|}{事件}&\multicolumn{1}{c}{经文}\\\hline

罗得家谱&11:10-32\\\hline
随亞伯蘭：吾珥-哈兰-迦南&12:1-5\\\hline
亞伯蘭在示劍/伯特利两次築壇&12:6-9\\\hline
因饑荒随亞伯蘭從南地下埃及，從南地去伯特利&12:10-20\\\hline
在伯特利和亞伯蘭牧人相争，與亞伯蘭分開，迁往所多玛&13:1-13\\\hline
被四王掳走，亞伯蘭從希伯崙北上到大馬士革左邊的何把，救羅得&14:1-16\\\hline
亞伯拉罕6次為羅得代求（50->45->40->30->20->10）&18:16-33\\\hline
神滅所多瑪與蛾摩拉,天使救罗得逃离&19:1-29\\\hline
摩押人亞捫人的起源&19:30-38\\\hline
\end{tabular}
\end{table}

\begin{table}[H]
\centering
\begin{tabular}{p{1.cm}|p{4.25cm}|p{4.25cm}|p{4.25cm}}\hline\hline
\multicolumn{1}{c|}{次数}&\multicolumn{1}{c|}{罗得}&\multicolumn{1}{c|}{亚伯兰}&\multicolumn{1}{c}{结局}\\\hline
一次&罗得被地上四王掳&亚伯兰率领318人救罗得&被神的祭司麦基洗得祝福\\\hline
二次&所多玛要被天上的神毁灭&亚伯兰在神前代求&为地上的亚比米勒祝福\\\hline
\end{tabular}
\end{table}



\subsubsection{新约对羅得的评价}

\subsection{使女之子以實瑪利——人意生的 husband's will}
\subsubsection{以實瑪利後代 25:12-18 }
\begin{forest}
  for tree={
    child anchor=west,
    parent anchor=east,
    grow=east,
    draw,
    anchor=west,
    edge path={
      \noexpand\path[\forestoption{edge}]
        (.child anchor) -| +(-5pt,0) -- +(-5pt,0) |-
        (!u.parent anchor)\forestoption{edge label};
    },
  }
  [亞伯蘭[夏甲（妾，使女）[以實瑪利[尼拜約][基達][亞德別][米比衫][米施瑪][度瑪][瑪撒]
												[哈大][提瑪][伊突][拿非施][基底瑪]]]
  ]
\end{forest}

\subsubsection{以實瑪利生平}
\begin{table}[H]
\begin{threeparttable}
\begin{tabular}{c c c}
\toprule
 \multicolumn{1}{L{\dimexpr0.5\linewidth-1\tabcolsep\relax}}{\bfseries 事件} 
&\multicolumn{1}{L{\dimexpr0.2\linewidth-1\tabcolsep\relax}}{\bfseries  备注}   
&\multicolumn{1}{L{\dimexpr0.1\linewidth-1\tabcolsep\relax}}{\bfseries 圣经经文} \\\hline

\multicolumn{1}{L{\dimexpr0.5\linewidth-1\tabcolsep\relax}}{\bfseries 下埃及得使女夏甲}  
& \multicolumn{1}{L{\dimexpr0.2\linewidth-1\tabcolsep\relax}}{\bfseries  见Section~\ref{sec:XiaAiJiDeXiaJia}}
& \multicolumn{1}{L{\dimexpr0.1\linewidth-1\tabcolsep\relax}}{\bfseries 12:10-20}\\ 

\multicolumn{1}{L{\dimexpr0.5\linewidth-1\tabcolsep\relax}}{\bfseries 撒萊將夏甲給亞伯蘭為妾}  
& \multicolumn{1}{L{\dimexpr0.2\linewidth-1\tabcolsep\relax}}{\bfseries  见Section~\ref{sec:XiaJiaYuGod}}
& \multicolumn{1}{L{\dimexpr0.1\linewidth-1\tabcolsep\relax}}{\bfseries 16:1-6}\\ 

\multicolumn{1}{L{\dimexpr0.5\linewidth-1\tabcolsep\relax}}{\bfseries 夏甲在曠野遇見神}  
& \multicolumn{1}{L{\dimexpr0.2\linewidth-1\tabcolsep\relax}}{\bfseries  见表~ \ref{table:GodYuXiaJia}}
& \multicolumn{1}{L{\dimexpr0.1\linewidth-1\tabcolsep\relax}}{\bfseries 16:7-14}\\ 

\multicolumn{1}{L{\dimexpr0.5\linewidth-1\tabcolsep\relax}}{\bfseries 以實瑪利出生}  
& \multicolumn{1}{L{\dimexpr0.2\linewidth-1\tabcolsep\relax}}{\bfseries  见Section~\ref{sec:YiShiMaLiChuSheng}}
& \multicolumn{1}{L{\dimexpr0.1\linewidth-1\tabcolsep\relax}}{\bfseries 16:15-16}\\ 

\multicolumn{1}{L{\dimexpr0.5\linewidth-1\tabcolsep\relax}}{\bfseries 以實瑪利被逐}  
& \multicolumn{1}{L{\dimexpr0.2\linewidth-1\tabcolsep\relax}}{\bfseries 见Section~\ref{sec:YiShiMaLiBeiZhu}}
& \multicolumn{1}{L{\dimexpr0.1\linewidth-1\tabcolsep\relax}}{\bfseries 21:8-21}\\ 
\multicolumn{1}{L{\dimexpr0.5\linewidth-1\tabcolsep\relax}}{\bfseries 以實瑪利後代}  
& \multicolumn{1}{L{\dimexpr0.2\linewidth-1\tabcolsep\relax}}{\bfseries 见Section~\ref{sec:YiShiMaLiHouDai}}
& \multicolumn{1}{L{\dimexpr0.1\linewidth-1\tabcolsep\relax}}{\bfseries 25:12-18}\\ \hline
\end{tabular}
\caption{以實瑪利生平}
\label{table:YiShiMaLiShengPing}
\end{threeparttable}
\end{table}



\subsubsection{新约对以實瑪利的评价}




\subsection{亞伯蘭之子以撒 17:1-18:15, 21:1-8}
\subsubsection{以撒生平}
\begin{table}[H]
\begin{threeparttable}
\begin{tabular}{c c c}
\toprule
 \multicolumn{1}{L{\dimexpr0.5\linewidth-1\tabcolsep\relax}}{\bfseries 事件} 
&\multicolumn{1}{L{\dimexpr0.2\linewidth-1\tabcolsep\relax}}{\bfseries  备注}   
&\multicolumn{1}{L{\dimexpr0.1\linewidth-1\tabcolsep\relax}}{\bfseries 圣经经文} \\\hline

\multicolumn{1}{L{\dimexpr0.5\linewidth-1\tabcolsep\relax}}{\bfseries 神应许肉身生子}  
& \multicolumn{1}{L{\dimexpr0.2\linewidth-1\tabcolsep\relax}}{\bfseries 见Section~\ref{sec:YaBoLanYuGodLiYue}}
& \multicolumn{1}{L{\dimexpr0.1\linewidth-1\tabcolsep\relax}}{\bfseries 15}\\ 

\multicolumn{1}{L{\dimexpr0.5\linewidth-1\tabcolsep\relax}}{\bfseries 割禮同天神應許以撒出生}  
& \multicolumn{1}{L{\dimexpr0.2\linewidth-1\tabcolsep\relax}}{\bfseries 见Section~\ref{sec:GodPromisYBLH}}
& \multicolumn{1}{L{\dimexpr0.1\linewidth-1\tabcolsep\relax}}{\bfseries  17:15-21 }\\ 

\multicolumn{1}{L{\dimexpr0.5\linewidth-1\tabcolsep\relax}}{\bfseries 神向撒拉應許生以撒  见}  
& \multicolumn{1}{L{\dimexpr0.2\linewidth-1\tabcolsep\relax}}{\bfseries 见Section~\ref{sec:GodPromiseIsaac}}
& \multicolumn{1}{L{\dimexpr0.1\linewidth-1\tabcolsep\relax}}{\bfseries 18:1-15}\\ 

\multicolumn{1}{L{\dimexpr0.5\linewidth-1\tabcolsep\relax}}{\bfseries 以撒出生 }  
& \multicolumn{1}{L{\dimexpr0.2\linewidth-1\tabcolsep\relax}}{\bfseries }
& \multicolumn{1}{L{\dimexpr0.1\linewidth-1\tabcolsep\relax}}{\bfseries 21:1-7}\\ 

\multicolumn{1}{L{\dimexpr0.5\linewidth-1\tabcolsep\relax}}{\bfseries 亞伯拉罕獻以撒 }  
& \multicolumn{1}{L{\dimexpr0.2\linewidth-1\tabcolsep\relax}}{\bfseries }
& \multicolumn{1}{L{\dimexpr0.1\linewidth-1\tabcolsep\relax}}{\bfseries 22:1-19}\\ 

\multicolumn{1}{L{\dimexpr0.5\linewidth-1\tabcolsep\relax}}{\bfseries 以撒母亲撒拉壽終，被埋葬}  
& \multicolumn{1}{L{\dimexpr0.2\linewidth-1\tabcolsep\relax}}{\bfseries 见表~\ref{table:QiuDi}}
& \multicolumn{1}{L{\dimexpr0.1\linewidth-1\tabcolsep\relax}}{\bfseries 23}\\ 

\multicolumn{1}{L{\dimexpr0.5\linewidth-1\tabcolsep\relax}}{\bfseries 以撒娶親 }  
& \multicolumn{1}{L{\dimexpr0.2\linewidth-1\tabcolsep\relax}}{\bfseries }
& \multicolumn{1}{L{\dimexpr0.1\linewidth-1\tabcolsep\relax}}{\bfseries 24}\\ 

\multicolumn{1}{L{\dimexpr0.5\linewidth-1\tabcolsep\relax}}{\bfseries 以撒兒子們（雅各以掃）出生 }  
& \multicolumn{1}{L{\dimexpr0.2\linewidth-1\tabcolsep\relax}}{\bfseries }
& \multicolumn{1}{L{\dimexpr0.1\linewidth-1\tabcolsep\relax}}{\bfseries 25:19-26}\\ 

\multicolumn{1}{L{\dimexpr0.5\linewidth-1\tabcolsep\relax}}{\bfseries 亞伯拉罕離世，以撒以实玛利埋葬父亲}  
& \multicolumn{1}{L{\dimexpr0.2\linewidth-1\tabcolsep\relax}}{\bfseries }
& \multicolumn{1}{L{\dimexpr0.1\linewidth-1\tabcolsep\relax}}{\bfseries 25:1-11}\\ 

\multicolumn{1}{L{\dimexpr0.5\linewidth-1\tabcolsep\relax}}{\bfseries 雅各用红豆汤騙長子名分 }  
& \multicolumn{1}{L{\dimexpr0.2\linewidth-1\tabcolsep\relax}}{\bfseries 见Section~\ref{sec:YaGePianZhangZiMingFen}}
& \multicolumn{1}{L{\dimexpr0.1\linewidth-1\tabcolsep\relax}}{\bfseries 25:27-34}\\ 

\multicolumn{1}{L{\dimexpr0.5\linewidth-1\tabcolsep\relax}}{\bfseries 以撒在基拉耳寄居 26:1-5 }  
& \multicolumn{1}{L{\dimexpr0.2\linewidth-1\tabcolsep\relax}}{\bfseries }
& \multicolumn{1}{L{\dimexpr0.1\linewidth-1\tabcolsep\relax}}{\bfseries }\\ 

\multicolumn{1}{L{\dimexpr0.5\linewidth-1\tabcolsep\relax}}{\bfseries 亞比米勒貪戀利百加 26:6-11}  
& \multicolumn{1}{L{\dimexpr0.2\linewidth-1\tabcolsep\relax}}{\bfseries }
& \multicolumn{1}{L{\dimexpr0.1\linewidth-1\tabcolsep\relax}}{\bfseries }\\ 

\multicolumn{1}{L{\dimexpr0.5\linewidth-1\tabcolsep\relax}}{\bfseries 亞比米勒驅逐以撒 26:12-16}  
& \multicolumn{1}{L{\dimexpr0.2\linewidth-1\tabcolsep\relax}}{\bfseries }
& \multicolumn{1}{L{\dimexpr0.1\linewidth-1\tabcolsep\relax}}{\bfseries }\\ 

\multicolumn{1}{L{\dimexpr0.5\linewidth-1\tabcolsep\relax}}{\bfseries 以撒挖井 26:17-25}  
& \multicolumn{1}{L{\dimexpr0.2\linewidth-1\tabcolsep\relax}}{\bfseries }
& \multicolumn{1}{L{\dimexpr0.1\linewidth-1\tabcolsep\relax}}{\bfseries }\\ 

\multicolumn{1}{L{\dimexpr0.5\linewidth-1\tabcolsep\relax}}{\bfseries 亞比米勒找以撒立約 26:26-33}  
& \multicolumn{1}{L{\dimexpr0.2\linewidth-1\tabcolsep\relax}}{\bfseries }
& \multicolumn{1}{L{\dimexpr0.1\linewidth-1\tabcolsep\relax}}{\bfseries }\\ 

\multicolumn{1}{L{\dimexpr0.5\linewidth-1\tabcolsep\relax}}{\bfseries 以掃娶二妻 26:34-35 见Section~\ref{sec:YiSauErQi}}  
& \multicolumn{1}{L{\dimexpr0.2\linewidth-1\tabcolsep\relax}}{\bfseries }
& \multicolumn{1}{L{\dimexpr0.1\linewidth-1\tabcolsep\relax}}{\bfseries }\\ 


\multicolumn{1}{L{\dimexpr0.5\linewidth-1\tabcolsep\relax}}{\bfseries 以撒給兒子們祝福 27:1-46}  
& \multicolumn{1}{L{\dimexpr0.2\linewidth-1\tabcolsep\relax}}{\bfseries }
& \multicolumn{1}{L{\dimexpr0.1\linewidth-1\tabcolsep\relax}}{\bfseries }\\ 

\multicolumn{1}{L{\dimexpr0.5\linewidth-1\tabcolsep\relax}}{\bfseries 以撒要祝福以掃 27:1-5b}  
& \multicolumn{1}{L{\dimexpr0.2\linewidth-1\tabcolsep\relax}}{\bfseries }
& \multicolumn{1}{L{\dimexpr0.1\linewidth-1\tabcolsep\relax}}{\bfseries }\\ 

\multicolumn{1}{L{\dimexpr0.5\linewidth-1\tabcolsep\relax}}{\bfseries 利百加怂恿雅各骗祝福 27:5a-17}  
& \multicolumn{1}{L{\dimexpr0.2\linewidth-1\tabcolsep\relax}}{\bfseries }
& \multicolumn{1}{L{\dimexpr0.1\linewidth-1\tabcolsep\relax}}{\bfseries }\\ 

\multicolumn{1}{L{\dimexpr0.5\linewidth-1\tabcolsep\relax}}{\bfseries 雅各照计欺父骗祝福 27:18-29}  
& \multicolumn{1}{L{\dimexpr0.2\linewidth-1\tabcolsep\relax}}{\bfseries }
& \multicolumn{1}{L{\dimexpr0.1\linewidth-1\tabcolsep\relax}}{\bfseries }\\ 

\multicolumn{1}{L{\dimexpr0.5\linewidth-1\tabcolsep\relax}}{\bfseries 以掃痛哭求父祝福 27:30-40}  
& \multicolumn{1}{L{\dimexpr0.2\linewidth-1\tabcolsep\relax}}{\bfseries }
& \multicolumn{1}{L{\dimexpr0.1\linewidth-1\tabcolsep\relax}}{\bfseries }\\ 

\multicolumn{1}{L{\dimexpr0.5\linewidth-1\tabcolsep\relax}}{\bfseries 利百加提议雅各回娘家娶妻 27:41-46}  
& \multicolumn{1}{L{\dimexpr0.2\linewidth-1\tabcolsep\relax}}{\bfseries }
& \multicolumn{1}{L{\dimexpr0.1\linewidth-1\tabcolsep\relax}}{\bfseries }\\ 

\multicolumn{1}{L{\dimexpr0.5\linewidth-1\tabcolsep\relax}}{\bfseries 以撒遣雅各於舅家娶妻 28:1-31:55}  
& \multicolumn{1}{L{\dimexpr0.2\linewidth-1\tabcolsep\relax}}{\bfseries }
& \multicolumn{1}{L{\dimexpr0.1\linewidth-1\tabcolsep\relax}}{\bfseries }\\ 

\multicolumn{1}{L{\dimexpr0.5\linewidth-1\tabcolsep\relax}}{\bfseries 以掃娶以實瑪利的女兒為妻 28:6-9 见Section~\ref{sec:YiSaiQuYiShiMaLiNuEr}}  
& \multicolumn{1}{L{\dimexpr0.2\linewidth-1\tabcolsep\relax}}{\bfseries }
& \multicolumn{1}{L{\dimexpr0.1\linewidth-1\tabcolsep\relax}}{\bfseries }\\ 

\multicolumn{1}{L{\dimexpr0.5\linewidth-1\tabcolsep\relax}}{\bfseries 以撒壽終 35:27-29}  
& \multicolumn{1}{L{\dimexpr0.2\linewidth-1\tabcolsep\relax}}{\bfseries }
& \multicolumn{1}{L{\dimexpr0.1\linewidth-1\tabcolsep\relax}}{\bfseries }\\ \hline

\end{tabular}
\caption{以撒生平}
\label{table:YiSaShengPing}
\end{threeparttable}
\end{table}


\subsubsection{以撒挖井 26:17-25 }

\begin{table}[!hbp]
\begin{threeparttable}
\begin{tabular}{@{}p{0.1\textwidth}*{4}{L{\dimexpr0.2\linewidth-8\tabcolsep\relax}}@{}}
\toprule
   地點 & 井名 & 原因              \\\hline
\multicolumn{1}{L{\dimexpr0.1\linewidth-2\tabcolsep\relax}}{\bfseries 基拉耳 }  
& \multicolumn{1}{L{\dimexpr0.1\linewidth-2\tabcolsep\relax}}{\bfseries 埃色}
& \multicolumn{1}{L{\dimexpr0.8\linewidth-2\tabcolsep\relax}}{\bfseries 以撒就離開那裡，在基拉耳谷支搭帳篷，住在那裡。當他父親亞伯拉罕在世之日所挖的水井，因非利士人在亞伯拉罕死後塞住了，以撒就重新挖出來，仍照他父親所叫的叫那些井的名字。以撒的僕人在谷中挖井，便得了一口活水井。基拉耳的牧人與以撒的牧人爭競說：這水是我們的．以撒就給那井起名叫埃色〔相爭〕因為他們和他相爭}\\\hline
\multicolumn{1}{L{\dimexpr0.1\linewidth-2\tabcolsep\relax}}{\bfseries 基拉耳 }  
& \multicolumn{1}{L{\dimexpr0.1\linewidth-2\tabcolsep\relax}}{\bfseries 西提拿 }
& \multicolumn{1}{L{\dimexpr0.8\linewidth-2\tabcolsep\relax}}{\bfseries 以撒的僕人又挖了一口井、他們又為這井爭競、因此以撒給這井起名叫西提拿。〔為敵〕 }\\\hline
\multicolumn{1}{L{\dimexpr0.1\linewidth-2\tabcolsep\relax}}{\bfseries 基拉耳 }  
& \multicolumn{1}{L{\dimexpr0.1\linewidth-2\tabcolsep\relax}}{\bfseries 利河伯 }
& \multicolumn{1}{L{\dimexpr0.8\linewidth-2\tabcolsep\relax}}{\bfseries 以撒離開那裡、又挖了一口井、他們不為這井爭競了、他就給那井起名叫利河伯．〔寬闊〕他說、耶和華現在給我們寬闊之地、我們必在這地昌盛。 }\\\hline
\multicolumn{1}{L{\dimexpr0.1\linewidth-2\tabcolsep\relax}}{\bfseries 別是巴 }  
& \multicolumn{1}{L{\dimexpr0.1\linewidth-2\tabcolsep\relax}}{\bfseries   }
& \multicolumn{1}{L{\dimexpr0.8\linewidth-2\tabcolsep\relax}}{\bfseries 以撒從那裡上別是巴去。當夜耶和華向他顯現、說、我是你父親亞伯拉罕的　神、不要懼怕、因為我與你同在、要賜福給你、並要為我僕人亞伯拉罕的緣故、使你的後裔繁多。以撒就在那裡築了一座壇、求告耶和華的名、並且支搭帳棚．他的僕人便在那裡挖了一口井。 }\\\hline
\end{tabular}
\caption{以撒挖井}
\label{table:IsaacWell}
\end{threeparttable}
\end{table}

\subsubsection{雅各以掃的祝福对比 27:27-29/27:38-40}
\begin{table}[H]
\begin{threeparttable}
\begin{tabular}{@{}p{0.1\textwidth}*{4}{L{\dimexpr0.2\linewidth-8\tabcolsep\relax}}@{}}
\toprule
   雅各 & 以掃            \\\hline
\multicolumn{1}{L{\dimexpr0.45\linewidth-2\tabcolsep\relax}}{\bfseries 我兒的香氣如同耶和華賜福之田地的香氣一樣。願　神賜你天上的甘露、地上的肥土、並許多五穀新酒．願多民事奉你、多國跪拜你．願你作你弟兄的主、你母親的兒子向你跪拜．凡咒詛你的、願他受咒詛．為你祝福的、願他蒙福。}  
& \multicolumn{1}{L{\dimexpr0.45\linewidth-2\tabcolsep\relax}}{\bfseries 地上的肥土必為你所住、天上的甘露必為你所得．你必倚靠刀劍度日、又必事奉你的兄弟、到你強盛的時候、必從你頸項上掙開他的軛。 （27:38-40）}\\\hline
\end{tabular}
\caption{以撒給兒子們的祝福}
\label{table:IsaacSons}
\end{threeparttable}
\end{table}
















\subsection{以掃——血气生的（by nature) 36 }
\subsubsection{以掃的俩个赫人妻子+以實瑪利的女兒 26:34-35/28:6-9}
\label{sec:YiSauErQi}



\subsubsection{以掃的子孫 36:1-19 }

\begin{forest}
  for tree={
    child anchor=west,
    parent anchor=east,
    grow=east,
    draw,
    anchor=west,
    edge path={
      \noexpand\path[\forestoption{edge}]
        (.child anchor) -| +(-5pt,0) -- +(-5pt,0) |-
        (!u.parent anchor)\forestoption{edge label};
    },
  }
  [以掃（以東）
  	[阿何利巴瑪（希未人祭便的孫女，亞拿的女兒）[耶烏施（族長）][雅蘭（族長）][可拉（族長）]]
  	[亞大（赫人以倫的女兒）[以利法[提幔（族長）][阿抹（族長）][洗玻（族長）][迦坦（族長）][基納斯（族長）][亞瑪力（族長）（以利法的妾，亭納所生]]]
  	[巴實抹（以實瑪利的女兒、尼拜約的妹子）[流珥[拿哈（族長）][謝拉（族長）][沙瑪（族長）][米撒（族長）]]]]
  ]
\end{forest}


\subsubsection{以東地的原住民 36:20-43}
\begin{itemize} 
\itemsep-.35em 
\item 那地原有的居民何利人西珥的子孫記在下面
\item 從何利人所出的族長，都在西珥地，按著宗族做族長
\end{itemize} 

\begin{forest}
  for tree={
    child anchor=west,
    parent anchor=east,
    grow=east,
    draw,
    anchor=west,
    edge path={
      \noexpand\path[\forestoption{edge}]
        (.child anchor) -| +(-5pt,0) -- +(-5pt,0) |-
        (!u.parent anchor)\forestoption{edge label};
    },
  }
  [何利人、西珥的子孫  	
  	[羅坍（族長）[何利][希幔]]	[亭納（羅坍的妹子）]	
  	[朔巴（族長）[亞勒文][瑪拿轄][以巴錄][示玻][阿南]]
  	[祭便（族長）[亞雅][亞拿（放他父親祭便的驢、遇著溫泉的、就是這亞拿）]]
  	[亞拿（族長）[底順][阿何利巴瑪（亞拿的女兒）]]
  	[底順（族長）[欣但][伊是班][益蘭][基蘭]]
  	[以察（族長）[辟罕][撒番][亞干]]
  	[底珊（族長）[烏斯][亞蘭]]]  
\end{forest}








\subsection{雅各 37-50}

\subsubsection{雅各生平}
\begin{table}[H]
\begin{threeparttable}
\begin{tabular}{c c c}
\toprule
 \multicolumn{1}{L{\dimexpr0.5\linewidth-1\tabcolsep\relax}}{\bfseries 事件} 
&\multicolumn{1}{L{\dimexpr0.2\linewidth-1\tabcolsep\relax}}{\bfseries  备注}   
&\multicolumn{1}{L{\dimexpr0.1\linewidth-1\tabcolsep\relax}}{\bfseries 圣经经文} \\\hline

\multicolumn{1}{L{\dimexpr0.5\linewidth-1\tabcolsep\relax}}{\bfseries }  
& \multicolumn{1}{L{\dimexpr0.2\linewidth-1\tabcolsep\relax}}{\bfseries }
& \multicolumn{1}{L{\dimexpr0.1\linewidth-1\tabcolsep\relax}}{\bfseries }\\ 

\multicolumn{1}{L{\dimexpr0.5\linewidth-1\tabcolsep\relax}}{\bfseries 以撒兒子們（雅各以掃）出生}  
& \multicolumn{1}{L{\dimexpr0.2\linewidth-1\tabcolsep\relax}}{\bfseries }
& \multicolumn{1}{L{\dimexpr0.1\linewidth-1\tabcolsep\relax}}{\bfseries 25:27-29}\\ 

\multicolumn{1}{L{\dimexpr0.5\linewidth-1\tabcolsep\relax}}{\bfseries }  
& \multicolumn{1}{L{\dimexpr0.2\linewidth-1\tabcolsep\relax}}{\bfseries }
& \multicolumn{1}{L{\dimexpr0.1\linewidth-1\tabcolsep\relax}}{\bfseries }\\ 

\multicolumn{1}{L{\dimexpr0.5\linewidth-1\tabcolsep\relax}}{\bfseries 雅各在伯特利（路斯）夢见神 28:10-22 见Section~\ref{sec:YaGeMengJianShen}}  
& \multicolumn{1}{L{\dimexpr0.2\linewidth-1\tabcolsep\relax}}{\bfseries }
& \multicolumn{1}{L{\dimexpr0.1\linewidth-1\tabcolsep\relax}}{\bfseries }\\ 

\multicolumn{1}{L{\dimexpr0.5\linewidth-1\tabcolsep\relax}}{\bfseries }  
& \multicolumn{1}{L{\dimexpr0.2\linewidth-1\tabcolsep\relax}}{\bfseries }
& \multicolumn{1}{L{\dimexpr0.1\linewidth-1\tabcolsep\relax}}{\bfseries }\\ 

\multicolumn{1}{L{\dimexpr0.5\linewidth-1\tabcolsep\relax}}{\bfseries }  
& \multicolumn{1}{L{\dimexpr0.2\linewidth-1\tabcolsep\relax}}{\bfseries }
& \multicolumn{1}{L{\dimexpr0.1\linewidth-1\tabcolsep\relax}}{\bfseries }\\ 

\multicolumn{1}{L{\dimexpr0.5\linewidth-1\tabcolsep\relax}}{\bfseries }  
& \multicolumn{1}{L{\dimexpr0.2\linewidth-1\tabcolsep\relax}}{\bfseries }
& \multicolumn{1}{L{\dimexpr0.1\linewidth-1\tabcolsep\relax}}{\bfseries }\\ 

\multicolumn{1}{L{\dimexpr0.5\linewidth-1\tabcolsep\relax}}{\bfseries }  
& \multicolumn{1}{L{\dimexpr0.2\linewidth-1\tabcolsep\relax}}{\bfseries }
& \multicolumn{1}{L{\dimexpr0.1\linewidth-1\tabcolsep\relax}}{\bfseries }\\ 

\multicolumn{1}{L{\dimexpr0.5\linewidth-1\tabcolsep\relax}}{\bfseries }  
& \multicolumn{1}{L{\dimexpr0.2\linewidth-1\tabcolsep\relax}}{\bfseries }
& \multicolumn{1}{L{\dimexpr0.1\linewidth-1\tabcolsep\relax}}{\bfseries }\\ 

\multicolumn{1}{L{\dimexpr0.5\linewidth-1\tabcolsep\relax}}{\bfseries }  
& \multicolumn{1}{L{\dimexpr0.2\linewidth-1\tabcolsep\relax}}{\bfseries }
& \multicolumn{1}{L{\dimexpr0.1\linewidth-1\tabcolsep\relax}}{\bfseries }\\ 

\multicolumn{1}{L{\dimexpr0.5\linewidth-1\tabcolsep\relax}}{\bfseries }  
& \multicolumn{1}{L{\dimexpr0.2\linewidth-1\tabcolsep\relax}}{\bfseries }
& \multicolumn{1}{L{\dimexpr0.1\linewidth-1\tabcolsep\relax}}{\bfseries }\\ 

\multicolumn{1}{L{\dimexpr0.5\linewidth-1\tabcolsep\relax}}{\bfseries }  
& \multicolumn{1}{L{\dimexpr0.2\linewidth-1\tabcolsep\relax}}{\bfseries }
& \multicolumn{1}{L{\dimexpr0.1\linewidth-1\tabcolsep\relax}}{\bfseries }\\ 

\multicolumn{1}{L{\dimexpr0.5\linewidth-1\tabcolsep\relax}}{\bfseries }  
& \multicolumn{1}{L{\dimexpr0.2\linewidth-1\tabcolsep\relax}}{\bfseries }
& \multicolumn{1}{L{\dimexpr0.1\linewidth-1\tabcolsep\relax}}{\bfseries }\\ 

\multicolumn{1}{L{\dimexpr0.5\linewidth-1\tabcolsep\relax}}{\bfseries }  
& \multicolumn{1}{L{\dimexpr0.2\linewidth-1\tabcolsep\relax}}{\bfseries }
& \multicolumn{1}{L{\dimexpr0.1\linewidth-1\tabcolsep\relax}}{\bfseries }\\ 

\multicolumn{1}{L{\dimexpr0.5\linewidth-1\tabcolsep\relax}}{\bfseries }  
& \multicolumn{1}{L{\dimexpr0.2\linewidth-1\tabcolsep\relax}}{\bfseries }
& \multicolumn{1}{L{\dimexpr0.1\linewidth-1\tabcolsep\relax}}{\bfseries }\\ 

\multicolumn{1}{L{\dimexpr0.5\linewidth-1\tabcolsep\relax}}{\bfseries }  
& \multicolumn{1}{L{\dimexpr0.2\linewidth-1\tabcolsep\relax}}{\bfseries }
& \multicolumn{1}{L{\dimexpr0.1\linewidth-1\tabcolsep\relax}}{\bfseries }\\ 

\multicolumn{1}{L{\dimexpr0.5\linewidth-1\tabcolsep\relax}}{\bfseries }  
& \multicolumn{1}{L{\dimexpr0.2\linewidth-1\tabcolsep\relax}}{\bfseries }
& \multicolumn{1}{L{\dimexpr0.1\linewidth-1\tabcolsep\relax}}{\bfseries }\\ 

\multicolumn{1}{L{\dimexpr0.5\linewidth-1\tabcolsep\relax}}{\bfseries }  
& \multicolumn{1}{L{\dimexpr0.2\linewidth-1\tabcolsep\relax}}{\bfseries }
& \multicolumn{1}{L{\dimexpr0.1\linewidth-1\tabcolsep\relax}}{\bfseries }\\ 
\end{tabular}
\caption{雅各生平}
\label{table:YaGeLife}
\end{threeparttable}
\end{table}


\subsubsection{雅各去世前的祝福 47:28-31，48，49:1-33 }

\begin{table}[!hbp]
\begin{threeparttable}
\begin{tabular}{@{}p{0.1\textwidth}*{4}{L{\dimexpr0.9\linewidth-8\tabcolsep\relax}}@{}}
\toprule
 \textbf{支派} &  \textbf{祝福}              \\\hline

\multicolumn{1}{L{\dimexpr0.1\linewidth-2\tabcolsep\relax}}{\bfseries 流便} 
&\multicolumn{1}{L{\dimexpr0.9\linewidth-2\tabcolsep\relax}}{\bfseries 流便哪、你是我的長子、是我力量強壯的時候生的、本當大有尊榮、權力超眾。但你放縱情慾、滾沸如水、必不得居首位、因為你上了你父親的床、污穢了我的榻。}  \\\hline

\multicolumn{1}{L{\dimexpr0.1\linewidth-2\tabcolsep\relax}}{\bfseries 西緬和利未} 
&\multicolumn{1}{L{\dimexpr0.9\linewidth-2\tabcolsep\relax}}{\bfseries 西緬和利未是弟兄、他們的刀劍是殘忍的器具。我的靈阿、不要與他們同謀、我的心哪、不要與他們聯絡、因為他們趁怒殺害人命、任意砍斷牛腿大筋．他們的怒氣暴烈可咒．他們的忿恨殘忍可詛．我要使他們分居在雅各家裡、散住在以色列地中。}  \\\hline

\multicolumn{1}{L{\dimexpr0.1\linewidth-2\tabcolsep\relax}}{\bfseries 猶大} 
&\multicolumn{1}{L{\dimexpr0.9\linewidth-2\tabcolsep\relax}}{\bfseries 猶大阿、你弟兄們必讚美你、你手必掐住仇敵的頸項、你父親的兒子們必向你下拜。猶大是個小獅子．我兒阿、你抓了食便上去．你屈下身去、臥如公獅、蹲如母獅、誰敢惹你。圭必不離猶大、杖必不離他兩腳之間、直等細羅〔就是賜平安者〕來到、萬民都必歸順。猶大把小驢拴在葡萄樹上、把驢駒拴在美好的葡萄樹上、他在葡萄酒中洗了衣服、在葡萄汁中洗了袍褂。他的眼睛必因酒紅潤、他的牙齒必因奶白亮。 }  \\\hline

\multicolumn{1}{L{\dimexpr0.1\linewidth-2\tabcolsep\relax}}{\bfseries 西布倫} 
&\multicolumn{1}{L{\dimexpr0.9\linewidth-2\tabcolsep\relax}}{\bfseries 西布倫必住在海口、必成為停船的海口．他的境界必延到西頓。}  \\\hline

\multicolumn{1}{L{\dimexpr0.1\linewidth-2\tabcolsep\relax}}{\bfseries 以薩迦} 
&\multicolumn{1}{L{\dimexpr0.9\linewidth-2\tabcolsep\relax}}{\bfseries 以薩迦是個強壯的驢、臥在羊圈之中．他以安靜為佳、以肥地為美、便低肩背重、成為服苦的僕人。  }  \\\hline

\multicolumn{1}{L{\dimexpr0.1\linewidth-2\tabcolsep\relax}}{\bfseries 但} 
&\multicolumn{1}{L{\dimexpr0.9\linewidth-2\tabcolsep\relax}}{\bfseries 但必判斷他的民、作以色列支派之一。但必作道上的蛇、路中的虺、咬傷馬蹄、使騎馬的墜落於後。耶和華阿、我向來等候你的救恩。}  \\\hline

\multicolumn{1}{L{\dimexpr0.1\linewidth-2\tabcolsep\relax}}{\bfseries 迦得} 
&\multicolumn{1}{L{\dimexpr0.9\linewidth-2\tabcolsep\relax}}{\bfseries 迦得必被敵軍追逼、他卻要追逼他們的腳跟。 }  \\\hline

\multicolumn{1}{L{\dimexpr0.1\linewidth-2\tabcolsep\relax}}{\bfseries 亞設} 
&\multicolumn{1}{L{\dimexpr0.9\linewidth-2\tabcolsep\relax}}{\bfseries 亞設之地必出肥美的糧食、且出君王的美味。 }  \\\hline

\multicolumn{1}{L{\dimexpr0.1\linewidth-2\tabcolsep\relax}}{\bfseries 拿弗他利} 
&\multicolumn{1}{L{\dimexpr0.9\linewidth-2\tabcolsep\relax}}{\bfseries 拿弗他利是被釋放的母鹿、他出嘉美的言語。}  \\\hline

\multicolumn{1}{L{\dimexpr0.1\linewidth-2\tabcolsep\relax}}{\bfseries 約瑟} 
&\multicolumn{1}{L{\dimexpr0.9\linewidth-2\tabcolsep\relax}}{\bfseries 約瑟是多結果子的樹枝、是泉旁多結果的枝子、他的枝條探出牆外。弓箭手將他苦害、向他射箭、逼迫他．但他的弓仍舊堅硬、他的手健壯敏捷．這是因以色列的牧者、以色列的磐石、就是雅各的大能者．你父親的　神、必幫助你、那全能者、必將天上所有的福、地裡所藏的福、以及生產乳養的福、都賜給你。你父親所祝的福、勝過我祖先所祝的福、如永世的山嶺、至極的邊界、這些福必降在約瑟的頭上、臨列那與弟兄迥別之人的頂上。    }  \\\hline

\multicolumn{1}{L{\dimexpr0.1\linewidth-2\tabcolsep\relax}}{\bfseries 便雅憫} 
&\multicolumn{1}{L{\dimexpr0.9\linewidth-2\tabcolsep\relax}}{\bfseries 便雅憫是個撕掠的狼、早晨要喫他所抓的、晚上要分他所奪的。}  \\\hline

\bottomrule
\end{tabular}
\caption{雅各給12支派的臨終祝福}
\label{table:JacobBlessing}
\end{threeparttable}
\end{table}

\subsection{約瑟 37-50}
\subsubsection{約瑟生平}



\subsubsection{以撒/雅各/約瑟的重要事件年龄推算过程}

\begin{tikzpicture}
\draw (-0.5,6.65) -- (-.5,22.5);   
 \draw (-0.5,23) node[below=3pt] {以撒} node[above=3pt] {};
 \draw (0.7,22.) node[below=3pt] {生雅各和以扫} node[above=3pt] {};
\draw (0.7,21.6) node[below=3pt] {（创25:24）} node[above=3pt] {};
 \draw (-1,22.) node[below=3pt] {60岁} node[above=3pt] {};
 \draw (-0.4,21.65) -- (-.6,21.65);  
 \draw  [dashed] (1.8,21.65) -- (4.6,21.65);  
\draw  [dashed] (5.4,21.65) -- (11.6,21.65);  
\draw  [dashed] (-0.1,20. ) -- (4.,20.); 
 \draw (-0.4,20.) -- (-.6,20.);  
\draw (-1.1,20.35) node[below=3pt] {100岁} node[above=3pt] {};
 \draw (-0.4,18.65) -- (-.6,18.65); 
\draw (1.37,19.) node[below=3pt] {给双子祝福（创27:4）} node[above=3pt] {};
 \draw [red, very thick, -](-0.4,17.65) -- (-.6,17.65); 
\draw (1.,18.) node[below=3pt] {\textcolor{red}{差雅各去舅家娶妻}} node[above=3pt] {};
\draw (0.7,17.5) node[below=3pt] {（创28:2）} node[above=3pt] {};
\draw (-1.0,18.) node[below=3pt] {\textcolor{red}{1?}岁} node[above=3pt] {};%
\draw (-0.4,6.65) -- (-.6,6.65);  
 \draw (0.3,7) node[below=3pt] {以撒过世} node[above=3pt] {};
\draw (-1.1,7) node[below=3pt] {180岁} node[above=3pt] {};

 \draw (5,2.5) -- (5,21.5);     
 \draw (5,22) node[below=3pt] {雅各} node[above=3pt] {};
\draw (4.9,20.6) -- (5.1,20.6);  
\draw (7.3,21.0) node[below=3pt] {雅各买长子名分（创25:33）} node[above=3pt] {};
\draw  [dashed] (9.5,20.6) -- (11.8,20.6); 
\draw (4.9,20.) -- (5.1,20.);  
\draw (4.58,20.35) node[below=3pt] {40岁} node[above=3pt] {}; 
\draw  [dashed] (5.5,20.) -- (10.8,20.); 
\draw (6.58,19.) node[below=3pt] {得祝福（创27:28）} node[above=3pt] {};
\draw (4.9,18.65) -- (5.1,18.65);  
\draw (6.70,18.) node[below=3pt] {\textcolor{red}{雅各去舅拉班家娶妻}} node[above=3pt] {};
\draw (4.50,18.) node[below=3pt] {\textcolor{red}{2?}岁} node[above=3pt] {};%
\draw  [red, very thick, -](4.9,17.65) -- (5.1,17.65); %
\draw (4.9,16.65) -- (5.1,16.65); 
\draw (7.70,17.) node[below=3pt] {娶利亚,7天后娶拉结（创29:28）} node[above=3pt] {};
\draw[>=triangle 45, <->] (4.9,16.65) -- (4.9,17.65);
\draw (4.50,17.5) node[below=3pt] {7年} node[above=3pt] {};
\draw (7.18,16.) node[below=3pt] {利亚生流便/西面/利未/犹大} node[above=3pt] {};
\draw (6.73,15.5) node[below=3pt] {辟拉亚生但/拿弗他利} node[above=3pt] {};
\draw (6.35,15.) node[below=3pt] {悉帕生迦得/亞設} node[above=3pt] {};
\draw (7.10,14.5) node[below=3pt] {利亚生以薩迦/西布倫/底拿} node[above=3pt] {};
\draw [red, very thick, -](4.9,13.65) -- (5.1,13.65); %
\draw (5.98,14.) node[below=3pt] {\textcolor{red}{拉结生約瑟}} node[above=3pt] {};
\draw (4.5,14.) node[below=3pt] {\textcolor{red}{3?}岁} node[above=3pt] {};%
\draw  [red, very thick, -](4.9,12.65) -- (5.1,12.65); %
\draw (4.5,13.) node[below=3pt] {\textcolor{red}{4?}岁} node[above=3pt] {};%
\draw (6.39,13.) node[below=3pt] {\textcolor{red}{雅各离开拉班家}} node[above=3pt] {};
\draw[>=triangle 45, <->] (4.9,12.65) -- (4.9,16.65);
\draw (4.30,15.5) node[below=3pt] {13年} node[above=3pt] {};
%\draw (4.10,15.1) node[below=3pt] {生11儿} node[above=3pt] {};
%\draw (4.40,14.7) node[below=3pt] {1女} node[above=3pt] {};
\draw (6.38,11.9) node[below=3pt] {底拿被示剑玷污} node[above=3pt] {};
\draw (6.,11) node[below=3pt] {便雅悯出生} node[above=3pt] {};
\draw (4.9,6.65) -- (5.1,6.65);
\draw (4.3,7.) node[below=3pt] {120岁} node[above=3pt] {};
\draw (6.6,7) node[below=3pt] {雅各埋葬父亲以撒} node[above=3pt] {};
\draw (7.4,4.) node[below=3pt] {\textcolor{blue}{雅各带全家下埃及（创47:9）}} node[above=3pt] {};
\draw  [dashed] (9.6,3.65) -- (10.7,3.65);
\draw (4.9,3.65) -- (5.1,3.65); 
\draw (4.3,4.) node[below=3pt] {\textcolor{blue}{130岁}} node[above=3pt] {};
 \draw (5.9,3) node[below=3pt] {雅各过世} node[above=3pt] {};
\draw (4.9,2.65) -- (5.1,2.65); 
\draw (4.3,3) node[below=3pt] {147岁} node[above=3pt] {};

\draw (12,0) -- (12,13.55); 
\draw (12,14.) node[below=3pt] {約瑟} node[above=3pt] {};
\draw  [dashed] (7.,13.65) -- (11.5,13.65);  
\draw (13.8,10.) node[below=3pt] {約瑟被兄长们卖埃及} node[above=3pt] {};
\draw (13.4,7.) node[below=3pt] {約瑟做埃及宰相} node[above=3pt] {};
\draw (11.9,6.65) -- (12.1,6.65);  
\draw (11.5,7.) node[below=3pt] {30岁} node[above=3pt] {};
\draw (13.6,6.) node[below=3pt] {埃及地开始大饥荒} node[above=3pt] {};
\draw (14.3,5.2) node[below=3pt] {約瑟兄长们二次下埃及买粮} node[above=3pt] {};
\draw (11.5,5.2) node[below=3pt] {39岁} node[above=3pt] {};
\draw (12.8,4.7) node[below=3pt] {（创45:6）} node[above=3pt] {};
\draw (11.9,5.65) -- (12.1,5.65); 
\draw (11.9,4.8) -- (12.1,4.8); 
\draw[>=triangle 45, <->] (11.9,4.8) -- (11.9,5.65);
\draw[>=triangle 45, <->] (11.9,6.65) -- (11.9,5.65);
\draw (11.5,6.5) node[below=3pt] {7年} node[above=3pt] {};
\draw (11.5,5.6) node[below=3pt] {2年} node[above=3pt] {};
\draw (11.2,4.) node[below=3pt] {\textcolor{blue}{$\approx$40岁}} node[above=3pt] {};
\draw (11.9,3.65) -- (12.1,3.65); 
\draw (13.8,4.) node[below=3pt] {\textcolor{blue}{約瑟在埃及迎接父亲}} node[above=3pt] {};

\draw (12,16) -- (12,21.5);     
 \draw (12,22) node[below=3pt] {以扫} node[above=3pt] {};
\draw (11.9,20.6) -- (12.1,20.6);  
\draw (14.2,21.0) node[below=3pt] {以扫卖长子名分（创25:33）} node[above=3pt] {};
\draw (11.9,20) -- (12.1,20.0);  
\draw (11.5,20.35) node[below=3pt] {40岁} node[above=3pt] {};
\draw (13.9,20.35) node[below=3pt] {以扫娶2妻（创26:34）} node[above=3pt] {};
\draw (11.9,18.65) -- (12.1,18.65);  
\draw (13.55,19.) node[below=3pt] {失祝福（创27:39）} node[above=3pt] {};
\draw (11.9,17.65) -- (12.1,17.65); 
\draw (14.14,18.) node[below=3pt] {娶以实马利女儿（创28:9）} node[above=3pt] {};




\draw (4.,1.) node[below=3pt] 
{\textcolor{red}{3? 雅各生约瑟的岁数}\textcolor{blue}{ = 雅各下埃及的岁数 - 约瑟在埃及迎接父亲的岁数}} node[above=3pt] {};

\draw (3.75,0.6) node[below=3pt] 
{\textcolor{blue}{ = 130岁 -  40岁  $\approx$ 90岁}} node[above=3pt] {};

\draw (3.8,0.2) node[below=3pt] {\textcolor{red}{4? 雅各离开拉班家的岁数}：\textcolor{blue}{雅各在拉班家住了20年（创31：38），}} node[above=3pt] {};

\draw (7.45,-0.2) node[below=3pt] {\textcolor{blue}{7年后结婚（创：29:28),结婚13年中依次生了12个儿子}} node[above=3pt] {};

\draw (7.6,-0.6) node[below=3pt] {\textcolor{blue}{和1个女儿。按照圣经记载12个孩子是一个接一个的出生。}} node[above=3pt] {};

\draw (6.8,-1.0) node[below=3pt] {\textcolor{blue}{所以雅各离开拉班家时约瑟<=1岁，雅各<=91岁}} node[above=3pt] {};

\draw (5.7,-1.4) node[below=3pt] {\textcolor{red}{2? 雅各去拉班家的岁数}：\textcolor{blue}{雅各在拉班家住了20年（创31：38），所以雅各去的时候<=71岁}} node[above=3pt] {};

\draw (3.28,-1.8) node[below=3pt] {\textcolor{red}{1? 以撒差雅各去拉班家的岁数}：\textcolor{blue}{= 71+60， 以撒在130岁左右}} node[above=3pt] {};

\end{tikzpicture}


\subsection{猶大}



\subsection{圣经中的女人}
\subsubsection{夏娃（女人） 16:1-6 }
\subsubsection{撒萊 16:1-6 }
\subsubsection{使女夏甲 12:10-20 }
\subsubsection{利百加 12:10-20 }
\subsubsection{利亚 12:10-20 }
\subsubsection{拉结 12:10-20 }
\subsubsection{他瑪 38 }



